% 中文毕业论文 LaTeX 模板
% 编译方式：请使用 XeLaTeX 或 LuaLaTeX 编译，不要使用 pdflatex

\documentclass[12pt,a4paper,openany,oneside]{ctexbook}

% 基础包
\usepackage{amsmath}
\usepackage{amssymb}
\usepackage{amsthm}
\usepackage{booktabs}
\usepackage{graphicx}
\graphicspath{{picture/}}
\usepackage{subfigure}
\usepackage[linesnumbered,ruled,vlined]{algorithm2e}
\usepackage{comment}
\usepackage{pgfplots}
\pgfplotsset{compat=1.18}
\usepackage{geometry}
\usepackage{hyperref}
\usepackage{cleveref}
\crefname{figure}{图}{图}
\crefname{equation}{公式}{公式}
\crefname{table}{表}{表}
\crefname{theorem}{定理}{定理}
\crefname{proposition}{命题}{命题}
\crefname{lemma}{引理}{引理}
\crefname{corollary}{推论}{推论}
\crefname{definition}{定义}{定义}
\renewcommand{\figurename}{图}
\renewcommand{\tablename}{表}

% 定理环境定义
\newtheorem{theorem}{定理}[section]
\newtheorem{proposition}[theorem]{命题}
\newtheorem{lemma}[theorem]{引理}
\newtheorem{corollary}[theorem]{推论}
\theoremstyle{definition}
\newtheorem{definition}[theorem]{定义}
\newtheorem{example}[theorem]{例}
\theoremstyle{remark}
\newtheorem{remark}[theorem]{注}
\renewcommand{\proofname}{证明}

% 页面设置
\geometry{left=3cm,right=2.5cm,top=2.5cm,bottom=2.5cm}
\linespread{1.5}  % 1.5倍行距

% 章节标题格式
\ctexset{
  chapter={
    format=\centering\zihao{3}\heiti,
    beforeskip=0.8cm,
    afterskip=0.5cm,
  },
  section={
    format=\zihao{4}\heiti,
  },
  subsection={
    format=\zihao{-4}\heiti,
  },
}

\begin{document}

% ==================== 封面 ====================
\begin{titlepage}
\centering
\vspace*{1.5cm}

{\zihao{-3}\heiti 分类号：}\underline{\makebox[3cm][c]{TP391}} \hfill {\zihao{-3}\heiti 学号：}\underline{\makebox[3cm][c]{}}

\vspace{2cm}

{\zihao{1}\songti\bfseries 硕士毕业论文}

\vspace{2cm}

{\zihao{2}\heiti 基于几何感知的\\多尺度状态空间点云分析方法}

\vspace{1cm}

{\zihao{3}\heiti Geometry-Aware Multi-Scale\\State Space Method for Point Cloud Analysis}

\vspace{3cm}

\begin{flushleft}
\zihao{-3}
\renewcommand\arraystretch{1.5}
\begin{tabular}{ll}
\heiti 作\quad 者： & \underline{\makebox[7cm][l]{张新元（Xinyuan Zhang）}} \\
\heiti 指导教师： & \underline{\makebox[7cm][l]{}} \\
\heiti 学科专业： & \underline{\makebox[7cm][l]{计算机科学与技术}} \\
\heiti 研究方向： & \underline{\makebox[7cm][l]{点云分析}} \\
\heiti 所在学院： & \underline{\makebox[7cm][l]{}} \\
\heiti 答辩日期： & \underline{\makebox[7cm][l]{}} \\
\end{tabular}
\end{flushleft}

\vfill

\end{titlepage}
\clearpage

% ==================== 中文摘要 ====================
\chapter*{摘\quad 要}
\addcontentsline{toc}{chapter}{摘要}

随着三维传感技术的快速发展，点云数据在自动驾驶、机器人感知与数字城市等领域得到广泛应用。近年来，状态空间模型（State Space Model，SSM）凭借线性复杂度下的长序列建模能力，被引入点云分析任务并取得一定进展。然而，对点云的序列化处理不可避免地破坏了点云在空间上的相邻性，导致在三维空间上相邻点在序列中被分离，削弱了模型对局部几何结构的感知能力。此外，现有SSM方法的状态转移矩阵在训练后保持固定，难以根据不同几何内容和尺度需求自适应调整，限制了层次化特征表示的构建。

针对上述问题，本文提出PointSS，一种基于几何感知的多尺度状态空间点云分析方法。为缓解序列化引发的空间邻域丢失问题，本文提出全局几何感知机制（GGAM）。GGAM通过窗口化邻域高效提取相对位置、法向量和曲率等几何先验，并结合双序列化并行建模与跨序列注意力机制，将几何信息显式注入到Mamba的状态传递过程中。

针对传统SSM状态转移参数与输入无关、难以适应多尺度建模需求的问题，本文设计了自适应尺度解耦状态空间模型（ASD-SSM）。ASD-SSM通过动态参数生成机制，依据各尺度的局部几何内容自适应生成状态转移参数，并利用尺度约束因子引导细尺度关注局部几何细节、粗尺度建模全局语义依赖，从而构建层次化的特征表示。

在空域建模的基础上，本文进一步提出基于Chebyshev多项式近似的多尺度图状态空间模型（Chebyshev多尺度SSM），从图多项式的视角与前述空域方法形成互补。该方法利用Chebyshev多项式对图拉普拉斯算子进行递推展开，将点云特征隐式地分解为多个空间尺度的分量，并为每个分量配备独立的Mamba模块进行差异化序列建模，通过可学习的权重自适应融合多尺度特征，使模型能够同时捕获局部细粒度结构与大范围上下文信息。

在S3DIS室内语义分割基准上，PointSS在不依赖预训练的条件下达到75.9\% mIoU，超越当前最优SSM方法PCM 6.1\%；在ModelNet40三维分类基准上达到94.6\%整体精度，超越所有SSM类对比方法。消融实验进一步验证了各模块设计的有效性。

\vspace{1em}
\noindent\textbf{关键词：}点云分析；状态空间模型；几何感知；多尺度特征学习；图多项式多尺度建模

% ==================== 英文摘要 ====================
\chapter*{Abstract}
\addcontentsline{toc}{chapter}{Abstract}

With the rapid development of 3D sensing technology, point cloud data has been widely applied in autonomous driving, robot perception, and digital city modeling. Recently, State Space Models (SSMs), with their long-sequence modeling capability at linear complexity, have been introduced to point cloud analysis and achieved promising results. However, serializing 3D point clouds into 1D sequences inevitably disrupts the spatial adjacency of point clouds—points that are adjacent in 3D space become separated in the sequence, weakening the model's ability to capture local geometric structures. Moreover, the state transition parameters in existing SSM-based methods remain input-invariant, preventing adaptive adjustment to the modeling requirements at different scales and limiting the construction of hierarchical feature representations.

To address these issues, this paper proposes PointSS, a geometry-aware multi-scale state space method for point cloud analysis. To mitigate the loss of spatial proximity caused by serialization, we propose a Global Geometry-Aware Mechanism (GGAM). GGAM efficiently extracts geometric priors—including relative positions, surface normals, and curvatures—via windowed neighborhoods. Combined with dual-serialization parallel modeling and cross-sequence attention, these geometric priors are explicitly injected into Mamba's state propagation.

To address the input-invariant state transition parameters in conventional SSMs that cannot adapt to multi-scale modeling requirements, we design an Adaptive Scale-Decoupled State Space Model (ASD-SSM). Through a dynamic parameter generation mechanism, ASD-SSM adaptively generates state transition parameters based on the local geometric content at each scale, and employs scale constraint factors to guide fine scales toward capturing local geometric details while coarse scales model global semantic dependencies, thereby constructing hierarchical feature representations.

Building upon spatial-domain modeling, this paper further proposes a Chebyshev polynomial approximation-based Multi-Scale Graph State Space Model (Chebyshev Multi-Scale SSM). This module complements the aforementioned spatial-domain methods from a graph polynomial perspective. By recursively expanding the graph Laplacian via Chebyshev polynomials, point cloud features are implicitly decomposed into multiple spatial-scale components, each equipped with an independent Mamba module for differentiated sequence modeling. Learnable fusion weights adaptively balance the contributions across scales, enabling the model to simultaneously capture fine-grained local structures and broad contextual information.

On the S3DIS indoor semantic segmentation benchmark, PointSS achieves 75.9\% mIoU without pre-training, surpassing the state-of-the-art SSM-based method PCM by 6.1\%. On the ModelNet40 3D classification benchmark, PointSS achieves 94.6\% overall accuracy, outperforming all SSM-based counterparts. Ablation studies further validate the effectiveness of each proposed module.

\vspace{1em}
\noindent\textbf{Keywords:} Point Cloud Analysis; State Space Models; Geometry-Aware; Multi-Scale Feature Learning; Graph Polynomial Multi-Scale Modeling

% ==================== 目录 ====================
\tableofcontents
\clearpage

%% main text
%%
%% Use \chapter commands to start a chapter
\chapter{绪论}
\label{chap:introduction}

\section{研究背景及意义}

随着三维传感技术的快速发展，点云作为三维空间的主要离散化表示形式，在自动驾驶、机器人感知、增强现实与数字城市建模等领域得到了广泛应用\cite{Vertex}。相较于二维图像，点云中每个点携带精确的三维空间坐标，能够直接表达物体的空间位置与表面几何形状，且不受光照、色彩等外观因素干扰，在复杂环境下的三维感知中具有独特优势。随着激光雷达等传感器的普及，可采集的点云规模与密度持续增长，如何从大规模点云中高效、准确地提取几何与语义信息，已成为三维场景理解的核心问题。

然而，点云的高效分析面临三方面固有挑战。其一，点云的无序性和不规则性使得传统卷积操作难以直接适用，需要专门设计排列不变的特征提取架构。其二，点云在物体边界和复杂曲面区域分布稀疏，使得局部几何特征的提取尤为困难。其三，随着激光雷达扫描分辨率的提升，大规模场景中的点数可达百万量级，对模型计算复杂度提出了严苛要求。早期的PointNet系列方法\cite{pointnet,PointNet++}和Transformer类方法\cite{pt,ptv3}在建模能力上取得了显著进展，但受限于全局池化导致的信息损失和自注意力机制$O(N^2)$的计算开销，在大规模场景下难以兼顾精度与效率。基于状态空间模型（SSM）的方法\cite{Mamba,PointMamba,pcm}凭借线性时间复杂度下的长序列建模能力，为上述问题提供了新的解决思路。然而，将SSM适配到三维点云时，序列化不可避免地破坏点云的空间近邻关系，造成几何信息丢失；同时，现有方法在多尺度层次化建模方面能力不足，这些问题仍是制约其性能的关键瓶颈。

因此，研究面向大规模点云场景的高效几何感知与层次化特征学习方法，不仅具有重要的理论价值，也对推动自动驾驶、机器人感知等实际应用的发展具有直接意义。

\section{国内外研究现状}

近年来，点云分析领域从基于几何先验的传统方法演进至深度学习驱动的端到端模型。本节从深度学习点云语义分析、几何特征增强、多尺度特征表达以及基于图多项式的分析方法四个维度梳理相关工作。

\subsection{基于深度学习的点云语义分析}
PointNet~\cite{pointnet}开创性地通过对称函数（最大池化）处理点云无序性问题，成为点云深度学习的基础框架，但其全局池化策略在局部几何细节建模方面存在明显不足。为克服这一局限，PointNet++~\cite{PointNetpp}引入分层采样策略，通过抽象层递进式地实现多尺度特征学习，显著增强了对局部结构的感知能力。后续工作进一步增强了局部特征建模能力：DGCNN~\cite{DGCNN}通过在特征空间中动态构建k近邻图结构，将每轮卷积的局部感受野自适应地调整为语义近邻而非几何近邻；此外亦有方法通过核点算子、残差MLP架构、位置自适应卷积核等策略进一步挖掘局部几何信息，从不同角度拓展了卷积式点云分析的设计空间。然而，这些方法的感受野受限于邻域半径，难以有效捕获点云中的长程依赖关系，这在大场景语义分割中表现得尤为明显。

Transformer通过自注意力机制为点云分析引入了全局建模能力，突破了局部感受野的限制。Point Transformer~\cite{pt}设计了向量注意力机制，将相对位置信息融入注意力计算，实现几何感知的特征聚合，在室内场景分割基准上取得了显著提升；PTv3~\cite{ptv3}进一步采用简化的序列化策略，通过基于序列化顺序的窗口注意力取代传统空间划分窗口注意力，并配合大规模预训练充分挖掘模型表示能力。尽管Transformer具有强大的全局建模能力，其$O(N^2)$的计算复杂度在大规模点云处理中成为显著瓶颈，迫使模型依赖粗粒度表示或滑动窗口策略，难以在保持计算效率的同时兼顾精细局部几何细节。

状态空间模型（SSM）~\cite{Mamba}以线性复杂度提供了有竞争力的替代方案。PointMamba~\cite{PointMamba}较早地将Mamba架构引入点云分析，通过空间填充曲线将三维点云序列化为一维序列以实现全局建模；PCM~\cite{pcm}引入层次化下采样以捕获多尺度特征。然而，将三维点云转换为一维点云序列时面临两个关键挑战：其一，序列化破坏了空间邻近点之间的邻域关系，导致局部几何结构的空间邻近性丢失；其二，统一的序列化策略难以同时捕获细粒度局部特征与长程上下文依赖，限制了多尺度表达能力。

\subsection{基于几何特征增强的点云分析方法}
显式建模曲面法向量、曲率等几何属性对准确理解三维结构至关重要。早期方法主要依赖原始三维坐标，但随着研究深入，越来越多工作开始引入更丰富的几何先验知识。PointGA~\cite{pointga}通过三角函数位置编码将原始三维坐标扩展为高维几何特征，并设计了位置差分自注意力机制以线性复杂度实现几何感知的特征表达。PointMSGT~\cite{pointmsgt}提出多尺度几何特征提取框架，其几何特征提取模块通过每个点的K个近邻，以相邻近邻对与中心点重构三角形结构，利用三角形质心、法向量和平面常数提取精细的局部几何关系，多尺度注意力模块则通过自适应地聚合不同尺度的特征，增强模型对复杂场景的理解能力。Geo-CNN~\cite{geo-cnn}设计了几何卷积操作，将边特征提取过程分解到三个正交基上，并根据边向量与基之间的角度聚合特征，从而在特征提取层次中保持欧几里得空间的几何结构。

除坐标扩展外，几何基元表示也是重要方向。RepSurf~\cite{repsurf}提出基于三角网格和伞状曲率的代表性曲面表示方法，通过对局部邻域构建伞形曲面结构提取法向量和曲率等几何特征，以增强对局部曲面结构的理解能力。PointVector~\cite{pointvector}引入向量表示捕获方向性几何信息，通过显式建模点之间的方向关系提升对边缘和曲面法向的感知能力。频域方法方面，PointWavelet~\cite{pointwavelet}在频谱域分析点云结构，通过小波变换实现多分辨率几何特征学习，为点云分析提供了更具表达力的特征空间。图方法在几何建模中也展现出优势：DGACN~\cite{dgacn}通过双图注意力卷积网络同时建模局部几何图和非局部语义图，捕获多层次的几何依赖关系；LGGCM~\cite{lggcm}提出局部-全局图卷积方法，在局部邻域内提取精细几何特征的同时建立全局点之间的长程依赖。针对大规模场景，LiDAR点云分割方法~\cite{lidar_segmentation}强调了结构化几何推理在处理稀疏、大范围点云时的重要性。然而，如何在保持计算效率的同时将显式几何先验与深度学习特征有效融合，并适应点云的稀疏性和非均匀分布，仍是一个开放性挑战。

\subsection{基于多尺度特征融合的点云分析方法}
多尺度特征融合自PointNet++~\cite{PointNetpp}引入分层集合抽象层以来被广泛采用，该方法通过渐进式下采样和局部特征聚合实现层次化特征学习。PointNeXt~\cite{pointnext}在改进训练与缩放策略的基础上重新审视PointNet++，验证了数据流水线层面的架构改进可带来显著增益。后续工作探索了多种融合策略：基于Transformer的方法如PTv2~\cite{ptv2}和PTv3~\cite{ptv3}采用具有层次化注意力机制的U型网络架构，在编码器-解码器结构中实现跨尺度特征传递；OctFormer~\cite{octformer2}扩展八叉树结构以实现高效的多尺度点云建模，通过树状层次结构自然地组织不同分辨率的特征；CCMNet~\cite{ccmnet}探索多尺度与跨类型特征交互，增强三维学习能力；RandLA-Net~\cite{Randla}验证了随机采样结合局部特征聚合模块可有效处理大规模场景，在保持计算效率的同时捕获多尺度信息。这些工作共同证明了多尺度建模在点云分析中的重要性。

近期基于SSM的方法尝试将线性复杂度引入点云分析。PointMamba~\cite{PointMamba}采用平坦非层次化的Mamba主干，利用空间填充曲线（Hilbert曲线及其转置变体）进行双序列化建模。该设计优先考虑简洁性和全局依赖的高效建模，但在单一特征粒度下处理所有点，缺乏先前工作证明有效的显式层次化结构。相比之下，PCM~\cite{pcm}引入层次化下采样以捕获多尺度特征，解决了PointMamba的单尺度局限。然而，PCM为每个尺度分配独立参数集，且各尺度内的状态转移矩阵$A$在训练后保持固定，所有Patch共享同一组$A$参数。步长参数$\Delta$虽使$\overline{A}$逐点变化，但无法根据不同Patch的几何内容自适应地调整状态转移特性。此外，为每个尺度分配独立参数导致参数量增长和内存开销增加。

\subsection{基于图多项式的点云分析方法}
基于Chebyshev多项式的图卷积方法为图神经网络提供了高效的谱域特征提取框架。Defferrard等人提出的ChebNet~\cite{chebnet}将谱域图卷积转化为Chebyshev多项式的递推展开：通过将归一化图拉普拉斯算子线性映射到$[-1,1]$区间，以$K$阶多项式近似谱滤波器，使图卷积在$K$跳邻域内局部化完成，计算复杂度与$K$成正比，避免了显式特征分解的高计算开销。Kipf和Welling提出的GCN~\cite{gcn}取$K=1$的一阶近似，以极简参数化在节点分类任务上验证了图多项式方法的有效性。两者共同表明，$K$阶Chebyshev展开天然对应$K$跳邻域的多尺度特征聚合：$K$较小时聚合局部细粒度结构，$K$增大时感受野逐步向外扩展，为多尺度特征学习提供了理论支撑。

将上述图谱方法引入点云分析，多项工作通过KNN构建局部图并以Chebyshev多项式进行特征提取。Wang等人~\cite{localspecgcn}在PointNet++框架下，对每个局部邻域构建$k$近邻图并以Chebyshev谱卷积替代逐点MLP，通过谱递归聚合策略证明了拓扑信息对局部特征学习的显著增益，实验表明谱图卷积能够联合建模邻域内所有点的关系，而逐点MLP的输出仅依赖单点自身信息。PointNGCNN~\cite{pointngcnn}进一步将Chebyshev图滤波器同时应用于坐标空间和特征空间，通过KNN邻域图的拉普拉斯矩阵进行递推卷积并以最大池化聚合中心点特征，在分类与分割任务上验证了图多项式滤波优于逐点MLP的结构特征提取能力。RGCNN~\cite{rgcnn}则在各层根据学习到的特征动态更新图拉普拉斯算子，并引入图信号平滑正则项约束训练过程，在点云分割任务上取得了良好的性能。

然而，上述方法均将Chebyshev多项式视为图卷积滤波器：$K$阶递推产生的各阶分量直接经线性加权求和后统一送入后续激活层，不同跳数的邻域信息在进入非线性变换之前已被混合，各阶分量共享同一套参数，模型无法为不同尺度的邻域信息学习独立的非线性变换。

本文提出的Chebyshev多尺度SSM改变了Chebyshev递推的使用方式：将其作为多尺度分解器而非融合滤波器，显式保留$K$个独立分量（0跳、1跳、2跳等），并为每个分量配备独立的Mamba模块进行序列建模，使模型能够针对不同空间尺度的邻域信息采用差异化的状态演化策略，在同一处理步骤内实现多尺度几何特征的分离建模与融合。这与ASD-SSM通过扩充窗口尺寸实现的跨窗口多尺度建模形成互补，两者共同构建了本文的两级多尺度框架。

\section{关键问题与挑战}
目前基于Mamba的点云分析方法面临两方面关键挑战：几何建模能力不足与多尺度建模能力不足。

几何建模能力不足的根源在于序列化操作与点云空间邻域定义之间存在结构性矛盾。当前方法通常通过空间填充曲线对点集进行序列化，在局部窗口内保持一定的空间邻近性。然而，序列化无法保证所有空间邻近点在序列中同样邻近，在三维空间中相距很近的点，可能在一维序列中相距很远，如\ref{fig:analysis}左侧图所示。Mamba 在顺序特征提取过程中依赖状态向量进行信息传递，状态会随新点的输入不断更新。当空间邻近的点在序列中相距甚远时，它们之间的信息在到达彼此隐藏状态之前往往已经衰减，导致模型难以捕捉局部几何结构，分割效果随之下降，如图\ref{fig:analysis}右侧图所示。这一问题在物体边界处尤为突出，以墙面与房柱的边界为例，房柱与墙面边界处的准确分割依赖于序列化后相距甚远的空间邻近点之间的交互，从而造成边界分割效果不佳。由于点云的法向量、曲率、边界等几何属性本质上由空间邻域点集共同决定，空间临近性的丢失会使模型在学习曲面过渡、物体边界与细长结构时出现信息缺失，从而影响复杂形状的精确建模与分割效果。

\begin{figure}[htbp]
	\centering
	\includegraphics[width=\linewidth]{analysis.JPG}
	\caption{序列化后空间临近性丢失可视化。（左）：点云序列化示意；（右）：分割预测结果}
	\label{fig:analysis}
\end{figure}

多尺度建模能力不足的问题源于两方面。其一，现有SSM方法的状态转移参数与输入无关：PointMamba采用非层次化的平坦结构，在单一特征粒度下处理所有点，缺乏层次化表达；PCM虽然引入了层次化下采样，但各尺度内的状态转移矩阵$A$在训练后保持固定，所有Patch共享同一组参数，步长参数$\Delta$虽使$\overline{A}$逐点变化，但无法根据不同Patch的几何内容和所处尺度自适应地调整状态转移特性，限制了多尺度特征表达能力。其二，上述空域方法的感受野相对固定，难以在同一模块中同时建模不同空间尺度的几何信息。点云特征本身兼具局部细粒度结构（如曲面细节、边缘走向）与大范围上下文结构（如物体整体形态、区域语义）两类信息，二者的建模需求本质上不同，尽管可以通过改变窗口大小实现多尺度序列建模，但不同尺度的特征在进入序列模型前已被混合，各尺度无法获得独立的非线性建模能力。如何在保持线性复杂度的前提下，将不同空间尺度的特征显式分离并针对各尺度采用差异化建模策略，仍是一个有待解决的问题。

\section{主要工作内容}

为解决上述问题，本文提出PointSS，一种基于几何感知的多尺度状态空间点云分析方法。PointSS从空域与图多项式两个互补视角出发，系统性地解决基于SSM的点云分析中几何信息丢失与多尺度建模能力不足的问题。如\cref{fig:architecture}为PointSS模型整体架构，图中括号内标注（点数，特征维度）表示各阶段的张量形状。其中数据处理部分包括数据增强与点云序列化，序列化编码方式采用PTv3\cite{ptv3}的方法。数据处理完成后，首先由全局几何感知机制（GGAM）对点云进行显式几何特征增强。特征学习部分采用基于PointNet的U型编码器-解码器架构,采用以窗口为单位对点云进行上采样和下采样的策略\cite{PointNet++,ptv3}。编码器与解码器内嵌入ASD-SSM和Chebyshev多尺度SSM进行多尺度点云建模。为确保每个窗口包含相同数量的点，部分点可能出现在多个窗口中（图中以橙色标注）。学习到的特征随后送入预测头，生成最终的分割或分类结果。相较于PCM\cite{pcm}，该方案仅在训练初始时执行一次序列化，后续通过对序列化点云进行池化实现下采样，而非在每个下采样阶段重新序列化，从而避免了重复的$\mathcal{O}(N\log N)$序列化开销。

\cref{fig:coder}为编码器的基本架构，包括池化层被用来进行下采样以及点云信息聚合，多个重复出现的ASD-SSM和Chebyshev多尺度SSM模块被用来进行深层次的特征理解。与编码器不同的是，解码器使用上采样来替换池化层，由于解码器也包含ASD-SSM和Chebyshev多尺度SSM，因此在上采样时能够学习到不同编码器得到的层次信息。

\begin{figure}[htbp]
	\centering
	\includegraphics[width=\linewidth]{architecture.jpg}
	\caption{PointSS整体架构}
	\label{fig:architecture}
\end{figure}

\begin{figure}[htbp]
	\centering
	\includegraphics[width=\linewidth]{coder.JPG}
	\caption{编码器架构}
	\label{fig:coder}
\end{figure}

针对序列化导致的几何信息丢失问题，本文提出全局几何感知机制（GGAM）。GGAM作为特征增强模块，通过窗口化邻域高效提取相对位置、法向量和曲率等几何先验，并结合双序列化并行建模与跨序列注意力机制，将几何信息显式注入到Mamba的状态传递过程中，有效补偿了几何信息丢失导致的几何学习困难。

针对传统SSM状态转移参数与输入无关、难以适应多尺度建模需求的问题，本文设计了自适应尺度解耦状态空间模型（ASD-SSM）。ASD-SSM通过动态参数生成机制，依据各尺度的局部几何内容自适应生成状态转移参数，并利用尺度约束因子引导细尺度关注局部几何细节、粗尺度建模全局语义依赖，从而构建层次化的特征表示。

在空域建模的基础上，本文进一步提出基于Chebyshev多项式近似的多尺度图状态空间模型（Chebyshev多尺度SSM），从图多项式的视角与前述空域方法形成互补。ASD-SSM主要在序列化后的空域进行多尺度建模，但序列化不可避免地引入了空间邻域信息的损失；Chebyshev多尺度SSM通过对图拉普拉斯算子进行Chebyshev多项式递推展开，将点云特征隐式地分解为多个空间尺度的分量，并为每个分量配备独立的Mamba模块进行差异化序列建模，通过可学习的权重自适应融合多尺度特征。这种设计使模型能够同时捕获局部细粒度结构与大范围上下文信息，与ASD-SSM共同构成两级多尺度框架。

综上，本文的主要贡献包括：

\begin{itemize}
	\item 提出全局几何感知机制（GGAM）。该机制通过窗口化邻域高效提取几何先验，并结合双序列化融合、跨序列注意力与自适应门控，将显式几何信息注入Mamba的状态传递过程，从而补偿序列化割裂导致的几何学习困难，提升对复杂边界与曲面结构的刻画能力。

	\item 提出自适应尺度解耦状态空间模型（ASD-SSM）。ASD-SSM通过动态参数生成机制，依据各尺度的局部几何内容自适应生成状态转移参数，并利用尺度约束因子使细尺度侧重捕获局部几何细节、粗尺度建模全局语义依赖，在保持参数高效的前提下实现层次化的多尺度特征学习。

	\item 提出基于Chebyshev多项式近似的多尺度图状态空间模型（Chebyshev多尺度SSM）。该方法利用Chebyshev多项式对图拉普拉斯算子进行递推展开，以$O(KN)$的线性复杂度将点云特征分解为多个空间尺度的分量，并为每个分量配备独立的Mamba模块进行差异化建模，与ASD-SSM从空域和图多项式两个正交维度共同构建两级多尺度框架。

	\item 在S3DIS和ModelNet40等标准基准上进行了广泛实验。在S3DIS室内语义分割基准上，PointSS在不依赖预训练的条件下达到75.9\% mIoU，超越当前最优SSM方法PCM 6.1\%；在ModelNet40三维分类基准上达到94.6\%整体精度，超越所有SSM类对比方法。详尽的消融实验验证了各组件设计的有效性。
\end{itemize}

\section{论文章节安排}

本文的组织结构如下：

第一章为绪论，介绍点云分析的研究背景及意义、国内外研究现状、本文的主要工作内容和论文的章节安排。

第二章介绍相关理论和技术，系统阐述点云数据表示、深度学习方法、Transformer模型、状态空间模型与Mamba架构、点云序列化技术等基础知识。

第三章介绍PointSS的整体架构设计，阐述U型编码器-解码器框架的设计动机、层次化编码器的下采样策略与特征学习块组织、解码器的上采样机制与跳跃连接设计，以及面向语义分割和点云分类的任务预测头设计。

第四章介绍全局几何感知机制（GGAM），详细阐述几何丢失问题的动机与分析、GGAM的方法设计以及消融实验验证。

第五章介绍多尺度点云状态空间建模，分别阐述ASD-SSM（基于序列化的多尺度建模）和Chebyshev多尺度SSM（基于图多项式展开的多尺度建模）的设计动机、方法细节和实验分析。

第六章为实验结果与分析，展示PointSS在S3DIS、ModelNet40等标准基准上的性能表现，并通过综合实验展示三个模块的协同作用。

第七章为总结与展望，总结本文的主要工作和贡献，并对未来的研究方向进行展望。






\chapter{相关理论和技术}


\section{点云及其表示}

点云是三维空间中点的集合，每个点通常包含位置坐标$(x, y, z)$以及可选的属性信息（如颜色、法向量、强度等）。点云数据可以通过多种方式获取，包括激光雷达（LiDAR）扫描、深度相机、多视角立体重建等。相比于传统的二维图像，点云能够直接表示三维物体的几何结构，在自动驾驶\cite{qi2018frustum}、机器人导航\cite{rusu20113d}、增强现实\cite{izadi2011kinectfusion}等领域具有重要应用价值。

点云数据具有以下几个显著特性：

\textbf{无序性（Permutation Invariance）}：点云是无序点集，改变点的排列顺序不应影响其表示的三维形状。这要求点云处理网络必须具有排列不变性或等变性，能够输出与输入点顺序无关的结果。

\textbf{稀疏性和不规则性}：与规则网格结构的图像不同，点云在三维空间中分布稀疏且不规则。这使得传统的卷积操作难以直接应用，需要设计专门的网络结构来处理这种不规则数据。

\textbf{局部结构重要性}：点云中点的局部邻域包含了重要的几何信息，如曲率、法向量等。有效提取和利用这些局部几何特征对于理解三维形状至关重要。

\textbf{尺度变化}：同一物体在不同距离或分辨率下采样得到的点云具有不同的点密度。鲁棒的点云处理方法需要对这种尺度变化具有适应性。

\textbf{表面分布性}：点云数据本质上是对物体表面的离散采样，所有点均分布于物体的几何表面之上，而非填充其内部体积。这种表面分布特性意味着点云仅能隐式地描述物体的外部几何形态，缺乏对内部结构的直接表达。因此，点云处理网络需要能够从离散的表面采样点中推断出连续的曲面信息，理解点与点之间所隐含的几何连接关系，并在仅有表面观测数据的条件下实现对三维形状的准确感知与理解。
\section{Transformer模型}

Transformer\cite{vaswani2017attention}最初是为自然语言处理任务设计的，其核心创新是自注意力机制（Self-Attention Mechanism），能够建模序列中任意两个位置之间的依赖关系。给定输入序列，Transformer通过查询（Query）、键（Key）、值（Value）三个线性变换计算注意力权重：

\begin{equation}
\text{Attention}(Q, K, V) = \text{softmax}\left(\frac{QK^T}{\sqrt{d_k}}\right)V
\end{equation}

其中$d_k$是键向量的维度。这种机制允许模型根据内容的相似性动态地聚合信息，而不依赖于位置的远近。

Transformer的优势在于能够并行处理整个序列，并且能够捕获长距离依赖关系。然而，自注意力机制的计算复杂度为$O(N^2)$，其中$N$是序列长度，这在处理大规模点云时会带来巨大的计算开销。

近年来，研究者开始将Transformer应用于点云分析。Point Transformer\cite{zhao2021point}将自注意力机制引入点云处理，通过在局部邻域内计算点间的注意力权重来聚合特征。Point Transformer V2\cite{wu2022point}进一步改进了网络架构，采用分组向量注意力（Grouped Vector Attention）和位置编码优化，提高了计算效率和性能。Point Transformer V3\cite{wu2024point}引入了序列化注意力机制，通过将点云序列化后应用窗口注意力，显著降低了计算复杂度。

尽管Transformer在点云分析中取得了优异性能，但其二次复杂度仍然限制了在大规模场景中的应用。这促使研究者探索更高效的序列建模方法。

\section{状态空间模型与Mamba}

\subsection{状态空间模型基础}

状态空间模型（State Space Models, SSMs）是控制理论中的经典工具，用于描述系统的动态行为。连续时间的状态空间模型可以表示为：

\begin{equation}
\begin{aligned}
\dot{h}(t) &= Ah(t) + Bx(t) \\
y(t) &= Ch(t) + Dx(t)
\end{aligned}
\end{equation}

其中$x(t) \in \mathbb{R}^L$是输入信号，$h(t) \in \mathbb{R}^N$是隐状态，$\dot{h}(t)$表示隐状态对时间的导数，$y(t) \in \mathbb{R}^L$是输出信号，$A \in \mathbb{R}^{N \times N}$是状态矩阵，$B \in \mathbb{R}^{N \times L}$、$C \in \mathbb{R}^{L \times N}$、$D \in \mathbb{R}^{L \times L}$是输入、输出和直通矩阵。

为了在离散时间序列上应用SSM，需要通过离散化技术将连续模型转换为离散形式。使用零阶保持（Zero-Order Hold, ZOH）方法，可得：

\begin{equation}
\begin{aligned}
h_t &= \bar{A}h_{t-1} + \bar{B}x_t \\
y_t &= Ch_t
\end{aligned}
\end{equation}

其中$\Delta$是步长参数，$\Delta A$和$\Delta B$分别表示步长与状态矩阵、输入矩阵的乘积，$\bar{A} = \exp(\Delta A)$，$\bar{B} = (\Delta A)^{-1}(\exp(\Delta A) - I) \cdot \Delta B$。

传统的SSM（如S4\cite{gu2022efficiently}）通常假设系统参数是时间不变的（Linear Time Invariance, LTI），即矩阵$A$、$B$、$C$在整个序列中保持不变。这种设计虽然能够通过卷积定理实现高效的并行计算，但限制了模型根据输入内容动态调整的能力。

\subsection{Mamba架构}

Mamba\cite{gu2024mamba}是一种选择性状态空间模型，其核心创新在于打破了LTI约束，使得SSM参数能够根据输入内容动态变化。具体而言，Mamba将$B$、$C$和步长$\Delta$设计为输入相关的函数：

\begin{equation}
\begin{aligned}
B_t &= \text{Linear}_B(x_t) \\
C_t &= \text{Linear}_C(x_t) \\
\Delta_t &= \text{softplus}(\text{Parameter} + \text{Broadcast}_D(\text{Linear}_1(x_t)))
\end{aligned}
\end{equation}

其中$\text{Linear}_B$和$\text{Linear}_C$分别为将输入投影到$N$维的线性层，$\text{Linear}_1$将输入投影到1维标量，$\text{Broadcast}_D$将该标量广播到$D$维（$D$为模型通道数），$\text{Parameter}$为可学习的偏置参数，$\text{softplus}$激活函数确保步长$\Delta_t$为正值。步长$\Delta_t$控制离散化的时间步幅，通过$\bar{A} = \exp(\Delta_t A)$直接影响状态衰减速率：较大的$\Delta_t$加速状态衰减，使模型更关注当前输入；较小的$\Delta_t$减缓衰减，使模型保留更多历史信息。这种选择性机制使得模型能够根据输入内容选择性地更新隐状态，增强了对重要信息的记忆能力和对不相关信息的过滤能力。

Mamba的计算过程可以概括为：首先将输入$x$通过线性层映射到高维空间，并分为两路进行门控；随后通过选择性SSM根据输入动态计算$B_t$、$C_t$、$\Delta_t$并更新隐状态$h_t$；之后使用SiLU激活函数对两路特征进行门控融合；最后通过线性层映射回原始维度。

Mamba的一个关键优势是线性时间复杂度$O(N)$。虽然选择性机制破坏了卷积形式，无法直接使用FFT加速，但Mamba通过硬件感知的并行扫描算法（Hardware-Aware Parallel Scan）在GPU上实现了高效计算。该算法利用现代GPU的层次化内存结构（HBM、SRAM），通过kernel融合和内存复用，最小化IO开销，实现了比Transformer更高的吞吐量。

Mamba在语言建模、序列分类等任务上展现出了优异性能，其线性复杂度使其能够处理超长序列。这些优势促使研究者探索将Mamba应用于点云分析。

\section{点云序列化}

将无序的点云数据转换为有序序列是应用序列建模方法（如Transformer、Mamba）的前提。点云序列化的核心挑战在于如何在保持计算效率的同时，尽可能保留点云的空间邻近关系和几何结构信息。

\textbf{Z-order曲线}（也称Morton曲线）是一种经典的空间填充曲线，通过交错点的坐标位与构建Morton码来确定排序。具体地，对于三维点$(x, y, z)$，将每个坐标的二进制表示进行交错排列，得到Morton码，然后按Morton码排序。Z-order曲线能够较好地保持局部空间邻近性，但在某些情况下会出现跳跃，导致空间上相邻的点在序列中距离较远。

\textbf{Hilbert曲线}是另一种空间填充曲线，相比Z-order曲线具有更好的局部性保持特性。Hilbert曲线通过递归分割空间并连接子空间的方式构建，使得曲线上相邻的点在空间中也更可能相邻。然而，Hilbert曲线的计算复杂度较高，在实时应用中可能带来额外开销。

Hilbert曲线具有良好的空间局部性保持特性\cite{lawder2000using}：序列中相邻的点在3D空间中也趋向于邻近。这一性质可以形式化表述为：对于序列位置相近的点$p_i$和$p_j$（即$|i-j|$较小），其欧氏距离$\|x_i - x_j\|$的期望较小。定义序列邻域$\mathcal{N}_{seq}(i, w) = \{p_{i-w/2}, \ldots, p_{i-1}, p_{i+1}, \ldots, p_{i+w/2}\}$为点$p_i$在序列上前后各$w/2$个点，Hilbert曲线的空间局部性保证了：
\begin{equation}
\mathbb{E}_{j \in \mathcal{N}_{seq}(i, w)}[\|x_i - x_j\|] \leq \mathbb{E}_{j \sim \text{Uniform}}[\|x_i - x_j\|]
\end{equation}
即序列邻居的平均距离显著小于随机采样点的距离。

基于这一性质，可以直接利用序列上的局部邻居构建邻域关系，避免传统$k$-近邻（KNN）搜索的$O(N\log N)$复杂度，将邻域构建降低至$O(N)$。

窗口化序列化策略将点云划分为多个局部窗口，在每个窗口内进行序列化，然后将各窗口序列拼接。这种方法能够在一定程度上缓解全局序列化带来的邻近性破坏问题，但窗口边界处的点仍然可能与其空间邻居分离。

现有的序列化方法都存在一定的局限性：完全的空间邻近关系保持与计算效率之间存在权衡。这正是本文提出全局几何感知机制的动机之一——通过显式提取和注入几何先验信息，补偿序列化过程中的几何信息损失。

\subsection{序列化下采样机制}

在基于序列化的点云分析方法中，下采样是构建多尺度特征表示的关键步骤。传统的点云下采样方法通常采用最远点采样（Farthest Point Sampling, FPS）或随机采样，但这些方法在序列化点云上应用时会破坏已建立的序列顺序，需要重新进行序列化操作。

**基于池化的序列化下采样**：为避免重复序列化的计算开销，本文采用基于池化的下采样策略。具体地，将序列化后的点云按固定窗口大小进行分组，每个窗口内通过池化操作（如最大池化或平均池化）生成代表性特征。设序列化点云为$Q = (q_1, q_2, \ldots, q_N)$，窗口大小为$W$，则第$i$个池化窗口为：
\begin{equation}
P_i = \{q_{(i-1) \cdot W + 1}, q_{(i-1) \cdot W + 2}, \ldots, q_{i \cdot W}\}
\end{equation}

**多尺度窗口扩展**：在多尺度架构中，不同层级采用不同的窗口扩展因子。设基础窗口大小为$P$，在尺度$s$下，实际窗口大小为$P \times F_s$，其中$F_s$为扩展因子。对于三尺度配置，通常设置$F_1 = 1$（细尺度），$F_2 = 2$（中尺度），$F_3 = 4$（粗尺度）。这种设计使得粗尺度能够聚合更大范围的上下文信息，而细尺度保持对局部细节的敏感性。

**边界处理与填充策略**：当序列长度不能被窗口大小整除时，最后一个不完整窗口需要特殊处理。本文采用重复填充策略：将窗口内最后一个有效点的特征重复填充至窗口大小。这种策略相比零填充能更好地保持特征的连续性，避免引入人工的不连续性。

**计算复杂度分析**：传统方法在每个下采样层都需要重新进行序列化，总复杂度为$O(L \cdot N \log N)$，其中$L$为网络层数。而基于池化的下采样方法仅需一次初始序列化，后续下采样的复杂度为$O(N)$，总复杂度降低至$O(N \log N + L \cdot N)$。在深层网络中，这种效率提升尤为显著。

\section{图信号处理与频谱分析}

\subsection{图信号处理基础}

点云可以自然地建模为图结构$\mathcal{G} = (\mathcal{V}, \mathcal{E})$，其中节点$\mathcal{V}$表示点，边$\mathcal{E}$表示点间的邻接关系（通常通过$k$-近邻或半径邻域定义）。每个点的特征$x_i$可视为图上的信号。图信号处理理论\cite{shuman2013emerging}将经典信号处理的概念推广到图域，其中图拉普拉斯算子扮演了傅里叶变换基的角色。

\textbf{图拉普拉斯算子}：定义为$L = D - W$，其中$D$为度矩阵（对角元素$D_{ii} = \sum_j W_{ij}$），$W$为邻接矩阵（$W_{ij}$表示点$i$和$j$之间的边权重）。归一化拉普拉斯定义为$L_{norm} = I - D^{-1/2}WD^{-1/2}$。

\textbf{频谱分解}：拉普拉斯算子$L$是实对称半正定矩阵，可以进行特征分解：
\begin{equation}
L = U\Lambda U^T
\end{equation}
其中$U = [u_0, u_1, \ldots, u_{N-1}] \in \mathbb{R}^{N \times N}$为特征向量矩阵（构成图傅里叶基），$\Lambda = \text{diag}(\lambda_0, \lambda_1, \ldots, \lambda_{N-1})$为特征值矩阵（频率）。图拉普拉斯的特征值对应频率：低特征值对应的特征向量在图上变化平滑，高特征值对应的特征向量变化剧烈\cite{shuman2013emerging}。在点云场景中，这种特性使得低频分量通常对应平坦区域（如墙面、地板），而高频分量对应几何突变区域（如物体边界、角点）。

\textbf{图傅里叶变换}：信号$x \in \mathbb{R}^N$的图傅里叶变换定义为：
\begin{equation}
\hat{x} = U^T x
\end{equation}
其中$\hat{x}_i = \langle x, u_i \rangle$表示信号在第$i$个频率上的分量。

\textbf{频谱滤波}：图卷积可以表示为频域中的滤波操作\cite{defferrard2016convolutional}：
\begin{equation}
y = g_\theta(L)x = Ug_\theta(\Lambda)U^Tx
\end{equation}
其中$g_\theta(\Lambda) = \text{diag}(\theta_0, \theta_1, \ldots, \theta_{N-1})$为频域滤波器。通过设计$\theta_i$的分布，可以实现低通滤波（保留低频）、高通滤波（保留高频）或带通滤波。

然而，直接进行特征分解的复杂度为$O(N^3)$，且需存储$N \times N$的特征向量矩阵，在大规模点云上不可行。Chebyshev多项式近似提供了一种高效的替代方案。

\subsection{Chebyshev多项式近似}

Chebyshev多项式近似\cite{defferrard2016convolutional,hammond2011wavelets}通过用低阶多项式逼近频域滤波器，避免显式计算特征分解。其理论基础是：任意平滑函数$g_\theta(\lambda)$可以在$[-1,1]$区间内用Chebyshev多项式展开：
\begin{equation}
g_\theta(\lambda) \approx \sum_{k=0}^{K-1} \theta_k T_k(\tilde{\lambda})
\end{equation}
其中$\theta_k$为第$k$阶的Chebyshev展开系数（在图神经网络中为可学习参数），$T_k$为$k$阶Chebyshev多项式，$\tilde{\lambda} = 2\lambda/\lambda_{max} - 1$为归一化的特征值（将$[0, \lambda_{max}]$映射到$[-1,1]$）。

Chebyshev多项式具有以下递推性质：
\begin{equation}
\begin{aligned}
T_0(x) &= 1 \\
T_1(x) &= x \\
T_k(x) &= 2xT_{k-1}(x) - T_{k-2}(x), \quad k \geq 2
\end{aligned}
\end{equation}

将该递推应用于拉普拉斯算子作用在信号$x$上，可得：
\begin{equation}
\begin{aligned}
T_0(\tilde{L})x &= x \\
T_1(\tilde{L})x &= \tilde{L}x = \left(\frac{2L}{\lambda_{max}} - I\right)x \\
T_k(\tilde{L})x &= 2\tilde{L}T_{k-1}(\tilde{L})x - T_{k-2}(\tilde{L})x
\end{aligned}
\end{equation}
其中$\tilde{L} = 2L/\lambda_{max} - I$为归一化拉普拉斯。

基于此递推，滤波操作可表示为：
\begin{equation}
y = g_\theta(L)x \approx \sum_{k=0}^{K-1} \theta_k T_k(\tilde{L})x
\end{equation}

关键优势在于，$T_k(\tilde{L})x$可通过递推计算，每步仅需$O(|\mathcal{E}|)$次稀疏矩阵-向量乘法（利用$L$的稀疏性），总复杂度为$O(K|\mathcal{E}|)$。对于$k$-近邻图，$|\mathcal{E}| = O(kN)$，故总复杂度为$O(KkN) = O(KN)$（$k$为常数）。这比直接特征分解的$O(N^3)$快数个数量级。

\textbf{多尺度解释}：不同阶的Chebyshev多项式$T_k(\tilde{L})$作用于图信号时，对应不同跳数的邻域聚合，与谱域中的频率成分存在隐式关联。具体而言，$T_0(\tilde{L})x = x$保留原始信号，即0跳的自身特征；$T_1(\tilde{L})x = \tilde{L}x$聚合1跳邻域的加权差异信息；更高阶的$T_k(\tilde{L})x$则对应$k$跳邻域聚合，感受野随$k$递增。

需要指出的是，$T_k(\tilde{L})$并非严格的带通滤波器，但低阶多项式更侧重于捕获平滑变化，高阶多项式能够表达更快速的空间变化\cite{defferrard2016convolutional}。通过选择合适的$K$（通常2-5），可以近似覆盖图信号的主要空间尺度。这一理论为本文多尺度图状态空间模型的设计奠定了基础。

\section{其他相关技术}

\subsection{多层感知机}

多层感知机（Multi-Layer Perceptron, MLP）是深度学习中最基础的网络结构，由多个全连接层堆叠而成。每个全连接层对输入进行线性变换，随后通过非线性激活函数引入非线性表达能力。在点云处理网络中，MLP广泛应用于特征变换、维度映射和分类头等模块。

MLP的基本构成包括线性变换和激活函数两部分。对于输入特征$x \in \mathbb{R}^{d_{in}}$，单层MLP的计算过程可表示为：
\begin{equation}
y = \sigma(Wx + b)
\end{equation}
其中$W \in \mathbb{R}^{d_{out} \times d_{in}}$为权重矩阵，$b \in \mathbb{R}^{d_{out}}$为偏置向量，$\sigma(\cdot)$为激活函数。线性变换$Wx + b$将输入特征从$d_{in}$维空间映射到$d_{out}$维空间，这一过程可以理解为对输入特征的加权组合和重组。激活函数则在线性变换之后引入非线性，使得网络能够拟合复杂的非线性关系。常用的激活函数包括ReLU、GELU和SiLU等，它们各有特点：ReLU计算简单但存在神经元死亡问题，GELU和SiLU则更加平滑，在Transformer和Mamba等现代架构中更为常用。

在本文的PointSS模型中，MLP主要用于以下几个场景。在特征提取阶段，MLP将原始点云坐标和属性映射到高维特征空间，为后续处理提供丰富的特征表示。在GGAM模块中，MLP用于生成查询、键、值向量，以及对几何特征进行变换。在Mamba模块中，MLP负责生成输入相关的SSM参数，如$B_t$、$C_t$和$\Delta_t$。在分类或分割任务的输出层，MLP将提取的特征映射到类别空间，完成最终的预测。

MLP的设计需要权衡表达能力和计算效率。增加层数和隐藏维度可以提升模型的拟合能力，但也会带来更多的参数量和计算开销。在实践中，通常采用两到三层的MLP，并通过残差连接和归一化技术来稳定训练。此外，MLP的权重初始化对训练效果有重要影响，常用的初始化方法包括Xavier初始化和Kaiming初始化，它们根据层的输入输出维度自适应地设置初始权重范围，避免梯度消失或爆炸。

\subsection{残差连接}

残差连接（Residual Connection）是现代深度神经网络的核心设计之一，最早由ResNet\cite{he2016deep}提出。其核心思想是在网络层之间引入跳跃连接，将输入直接加到输出上，使得网络学习的是残差映射而非完整映射。对于某个网络层或模块$\mathcal{F}(\cdot)$，残差连接的计算形式为：
\begin{equation}
y = \mathcal{F}(x) + x
\end{equation}
其中$x$为输入，$\mathcal{F}(x)$为该层学习的变换，$y$为最终输出。这一简单而有效的设计极大地缓解了深层网络训练中的梯度消失问题，使得训练数百层甚至上千层的网络成为可能。

深层神经网络在训练过程中面临梯度消失的困境。在反向传播过程中，梯度需要从输出层逐层传递到输入层，每经过一层都会与该层的权重矩阵相乘。当网络层数很深时，梯度在传播过程中会不断衰减，导致靠近输入层的参数几乎无法得到有效更新，网络难以收敛。残差连接通过引入恒等映射的捷径，为梯度提供了一条直通的传播路径。在反向传播时，损失对输入$x$的梯度为：
\begin{equation}
\frac{\partial \mathcal{L}}{\partial x} = \frac{\partial \mathcal{L}}{\partial y} \left(\frac{\partial \mathcal{F}(x)}{\partial x} + I\right)
\end{equation}
其中$I$为单位矩阵。$\mathcal{L}$表示损失函数。即使$\frac{\partial \mathcal{F}(x)}{\partial x}$很小，梯度仍然可以通过恒等映射$I$直接传递，从而保证了深层网络的可训练性。

残差连接的另一个重要作用是促进特征复用。在没有残差连接的网络中，每一层必须从头学习完整的特征变换。而有了残差连接后，每一层只需要学习在输入基础上的增量变化，这使得网络更容易优化。从优化的角度看，残差连接为网络提供了更平滑的损失曲面，降低了优化难度。实验表明，即使某些残差块学习到的变换接近零映射，网络仍然能够保持良好的性能，因为输入可以通过跳跃连接直接传递到后续层。

在本文的PointSS模型中，残差连接被广泛应用于各个模块。在Mamba块中，每个Mamba层的输出都会与输入相加，形成残差结构。在GGAM模块中，几何特征增强后的结果通过残差连接与原始特征融合。在整体网络架构中，每个阶段的输出也通过残差连接与输入特征结合，确保信息在网络中的顺畅流动。这些残差连接不仅稳定了训练过程，也提升了模型的最终性能。

需要注意的是，残差连接要求输入和输出的维度必须一致。当维度不匹配时，通常通过线性投影层对输入进行维度调整，使其与输出维度对齐后再进行相加。此外，残差连接通常与归一化技术配合使用，形成"归一化-变换-残差连接"的标准模式，这一模式在Transformer、Mamba等现代架构中已成为标配。

\subsection{交叉熵损失函数}

交叉熵损失函数（Cross-Entropy Loss）是分类任务中最常用的损失函数，用于衡量模型预测的概率分布与真实标签分布之间的差异。在点云语义分割任务中，交叉熵损失计算每个点的预测类别与真实类别之间的误差，并通过反向传播指导网络参数的更新。

对于$C$类分类问题，给定模型输出的logits向量$z \in \mathbb{R}^C$，首先通过softmax函数将其转换为概率分布：
\begin{equation}
p_i = \frac{e^{z_i}}{\sum_{j=1}^{C} e^{z_j}}, \quad i = 1, 2, \ldots, C
\end{equation}
其中$p_i$表示预测为第$i$类的概率。对于真实类别$y$，交叉熵损失定义为：
\begin{equation}
\mathcal{L}_{CE} = -\log(p_y) = -z_y + \log\left(\sum_{j=1}^{C} e^{z_j}\right)
\end{equation}
交叉熵损失衡量的是预测概率分布与真实分布的差异：当模型对正确类别的预测概率越高，损失越小；反之，如果模型对错误类别给出较高概率，损失会显著增大。这种非对称的惩罚机制促使模型不断提高对正确类别的置信度。

在实际应用中，交叉熵损失具有良好的数值稳定性和优化特性。由于softmax函数和对数函数的组合，交叉熵损失的梯度形式非常简洁，梯度大小与预测误差成正比，这使得模型在预测错误时能够获得较大的梯度信号，加速收敛。此外，交叉熵损失是凸函数，在单层网络中保证了全局最优解的存在，虽然深度网络的损失曲面是非凸的，但交叉熵的良好性质仍然有助于优化。

在点云语义分割任务中，交叉熵损失需要处理类别不平衡问题。真实场景中，不同类别的点数量往往差异巨大，例如地面和墙面的点数远多于小物体如椅子和桌子的点数。如果直接使用标准交叉熵，模型会倾向于预测多数类别，而忽略少数类别。为解决这一问题，常用的策略包括类别加权和Focal Loss。类别加权为每个类别赋予不同的权重，少数类别的权重更高，从而平衡不同类别对损失的贡献。Focal Loss则通过降低易分类样本的权重，使模型更关注难分类样本，进一步缓解类别不平衡问题。

本文在训练PointSS模型时采用了带类别权重的交叉熵损失。权重根据各类别在训练集中的点数反比例设置，确保模型对所有类别都能给予足够的关注。此外，在计算损失时会忽略标注为无效的点，避免噪声标签对训练的干扰。

\subsection{AdamW优化器}

优化器负责根据损失函数的梯度更新网络参数，是深度学习训练过程中的核心组件。AdamW\cite{loshchilov2018decoupled}是Adam优化器的改进版本，通过解耦权重衰减与梯度更新，在训练Transformer、Mamba等现代深度网络时表现出色，已成为当前最流行的优化器之一。

Adam（Adaptive Moment Estimation）优化器结合了动量法和自适应学习率两种思想。其参数更新规则为：
\begin{equation}
\begin{aligned}
m_t &= \beta_1 m_{t-1} + (1-\beta_1) g_t \\
v_t &= \beta_2 v_{t-1} + (1-\beta_2) g_t^2 \\
\hat{m}_t &= \frac{m_t}{1-\beta_1^t}, \quad \hat{v}_t = \frac{v_t}{1-\beta_2^t} \\
\theta_t &= \theta_{t-1} - \eta \frac{\hat{m}_t}{\sqrt{\hat{v}_t} + \epsilon}
\end{aligned}
\end{equation}
其中$g_t$为第$t$步的梯度，$m_t$和$v_t$分别为梯度的一阶矩（均值）和二阶矩（未中心化的方差）的指数移动平均，$\beta_1$和$\beta_2$为衰减率（通常取0.9和0.999），$\hat{m}_t$和$\hat{v}_t$为偏差修正后的矩估计，$\eta$为学习率，$\epsilon$为数值稳定项（通常为$10^{-8}$）。动量法通过累积历史梯度的指数移动平均来加速收敛，使得参数更新方向更加稳定，能够有效穿越损失曲面的平坦区域。自适应学习率则为每个参数维度维护独立的学习率，根据历史梯度的二阶矩（方差）自动调整步长：梯度变化剧烈的参数使用较小的学习率，梯度平稳的参数使用较大的学习率。

然而，Adam在处理权重衰减（L2正则化）时存在理论缺陷。传统的权重衰减是在损失函数中加入参数范数的惩罚项，其梯度会与数据梯度一起参与Adam的自适应学习率计算。这导致权重衰减的效果被自适应学习率稀释，无法达到预期的正则化效果。AdamW通过将权重衰减从梯度计算中解耦出来，直接在参数更新时施加衰减：
\begin{equation}
\theta_t = \theta_{t-1} - \eta \frac{\hat{m}_t}{\sqrt{\hat{v}_t} + \epsilon} - \eta \lambda \theta_{t-1}
\end{equation}
其中$\lambda$为权重衰减系数。这种解耦使得权重衰减的强度不受自适应学习率的影响，从而恢复了权重衰减的正则化能力。

AdamW的另一个优势是对超参数的鲁棒性。相比SGD等优化器，AdamW对学习率的选择不那么敏感，通常使用较小的学习率（如1e-4到1e-3）即可取得良好效果。权重衰减系数通常设置在0.01到0.1之间，根据模型大小和任务特点调整。此外，AdamW的两个动量参数通常使用默认值：一阶动量系数为0.9，二阶动量系数为0.999，这些默认值在大多数任务中都表现良好。

在训练PointSS模型时，本文采用AdamW优化器，学习率设置为0.002，权重衰减系数为0.05。这些超参数在ScanNet和S3DIS等数据集上经过验证，能够在保证收敛速度的同时避免过拟合。AdamW的自适应特性使得模型在训练初期能够快速下降，在后期则稳定收敛，无需复杂的学习率调整策略。

\subsection{OneCycleLR学习率调度策略}

学习率是影响深度学习训练效果的关键超参数。固定学习率往往难以在收敛速度和最终性能之间取得平衡：过大的学习率导致训练不稳定甚至发散，过小的学习率则收敛缓慢且容易陷入局部最优。学习率调度策略通过在训练过程中动态调整学习率，能够兼顾快速收敛和精细优化。OneCycleLR\cite{smith2019super}是一种高效的学习率调度策略，通过单周期的学习率变化实现快速训练和优异性能。

OneCycleLR的核心思想是在训练过程中让学习率经历一个完整的周期：从较小的初始值逐渐上升到最大值，然后再逐渐下降到远小于初始值的最终值。这一策略基于超收敛（Super-Convergence）现象的观察：在训练初期使用较大的学习率可以帮助模型快速探索参数空间，跳出尖锐的局部最优，找到更平坦、泛化性能更好的区域；在训练后期使用较小的学习率则能够在找到的良好区域内进行精细优化，提升最终性能。

OneCycleLR的调度过程通常分为三个阶段。在上升阶段，学习率从初始值线性或余弦上升到最大值，这一阶段通常占据总训练步数的30%到45%。较大的学习率使得模型能够快速学习数据的主要模式，同时起到正则化作用，防止模型过早地拟合训练数据的噪声。在下降阶段，学习率从最大值逐渐下降到初始值以下，这一阶段占据大部分训练时间。学习率的下降使得模型在已经找到的良好区域内进行精细调整，逐步收敛到最优解。在最终阶段，学习率降至极小值，模型进行最后的微调，确保收敛到稳定的参数配置。

OneCycleLR相比传统的阶梯式学习率衰减（Step Decay）和余弦退火（Cosine Annealing）具有明显优势。阶梯式衰减在学习率突然下降时可能导致训练不稳定，且需要手动设置衰减的时机和幅度。余弦退火虽然平滑，但通常需要多个周期才能达到最佳效果，训练时间较长。OneCycleLR通过单个周期即可完成训练，大幅缩短训练时间，同时由于学习率的动态变化起到了正则化作用，模型的泛化性能往往更好。

在实际应用中，OneCycleLR的关键超参数包括最大学习率、初始学习率与最大学习率的比值、以及上升阶段占总步数的比例。最大学习率通常通过学习率范围测试（Learning Rate Range Test）确定：在训练初期以指数增长的方式增加学习率，观察损失的变化，选择损失下降最快时的学习率作为最大学习率。初始学习率通常设置为最大学习率的1/10到1/25，最终学习率则设置为初始学习率的1/100到1/1000。上升阶段的比例通常设置为30%，这一比例在多数任务中表现良好。

本文在训练PointSS模型时采用OneCycleLR策略，最大学习率设置为0.002，初始学习率为0.0002，最终学习率为0.00002，上升阶段占总训练步数的30%。这一配置在ScanNet数据集上能够在50个epoch内达到收敛，相比传统的阶梯式衰减策略节省了约20%的训练时间，同时模型的分割精度也有小幅提升。OneCycleLR的自动化特性减少了手动调整学习率的工作量，使得训练过程更加高效和稳定。

\subsection{注意力机制变体}

为了降低标准自注意力的二次复杂度，研究者提出了多种高效注意力机制。线性注意力（Linear Attention）\cite{katharopoulos2020transformers}通过核技巧将softmax注意力近似为线性操作，将复杂度降至$O(N)$，但可能损失部分表达能力。FlashAttention\cite{dao2022flashattention}通过优化内存访问模式和kernel融合，在保持精确注意力计算的同时显著提升了效率。稀疏注意力（Sparse Attention）仅计算部分位置对之间的注意力，通过限制注意力范围来降低计算量。

\subsection{位置编码}

由于Transformer本身不包含位置信息，需要通过位置编码注入序列的顺序信息。绝对位置编码（如正弦位置编码\cite{vaswani2017attention}）为每个位置生成固定的编码向量。相对位置编码（如旋转位置编码RoPE\cite{su2024roformer}）则编码位置之间的相对关系，通常具有更好的泛化能力。在点云处理中，位置编码需要结合三维空间坐标的特性进行设计。

\subsection{归一化技术}

层归一化（LayerNorm）\cite{ba2016layer}和RMSNorm\cite{zhang2019root}是深度网络中常用的归一化技术。LayerNorm计算每个样本在特征维度上的均值和方差进行归一化，而RMSNorm仅使用均方根进行归一化，计算更高效。归一化技术能够稳定训练过程，加速收敛。

\subsection{激活函数}

除了传统的ReLU激活函数，现代深度学习网络常使用更平滑的激活函数。GELU（Gaussian Error Linear Unit）\cite{hendrycks2016gaussian}基于输入的累积分布函数进行激活，在Transformer中广泛应用。SiLU（Sigmoid Linear Unit，也称Swish）\cite{ramachandran2017searching}结合了sigmoid函数的平滑性和线性函数的简单性。GLU（Gated Linear Unit）\cite{dauphin2017language}及其变体（如SwiGLU）通过门控机制控制信息流动，在语言模型和视觉模型中都取得了良好效果。Mamba架构中使用了SiLU和门控机制来增强非线性表达能力。

\section{本章小结}

本章系统介绍了本文研究所涉及的相关理论和技术。首先介绍了点云数据的基本特性，包括无序性、稀疏性和局部结构重要性。然后概述了深度学习在点云处理中的应用，包括PointNet、PointNet++、DGCNN等经典方法。接着详细讨论了Transformer模型的自注意力机制及其在点云分析中的应用。随后深入介绍了状态空间模型的数学基础和Mamba架构的选择性机制，阐明了其线性复杂度的优势和硬件感知优化。点云序列化技术的讨论揭示了现有方法在空间邻近性保持方面的局限性，这正是本文提出全局几何感知机制的重要动机。此外，本章还系统介绍了图信号处理与频谱分析的理论基础，包括图拉普拉斯算子、频谱分解、图傅里叶变换等核心概念，以及Chebyshev多项式近似方法，为后续多尺度图状态空间模型的设计提供了理论支撑。最后介绍了注意力机制变体、位置编码、归一化技术、激活函数等相关技术。这些理论和技术为后续章节中PointSS方法的设计提供了坚实基础。


\chapter{PointSS整体架构设计}

绪论中已从宏观层面介绍了PointSS的整体思路。本章对PointSS的网络架构进行系统性阐述，重点说明各组件的具体设计方案与设计动机，为后续各创新模块章节的展开提供全局视角。

\section{整体框架}

PointSS采用U型编码器-解码器结构（U-Net），如\cref{fig:architecture}所示，整体前向流程分为输入处理、编码器、解码器和任务预测头四个阶段。输入点云经序列化预处理后，依次经过GGAM几何特征增强、多阶段编码器逐层下采样提取深层特征，再由解码器逐层上采样恢复空间分辨率，最终通过任务预测头输出分割或分类结果。

U型结构的选择出于以下考量：其一，编码器的逐层下采样使感受野随深度增大，低层关注细粒度局部几何，高层建模场景级语义上下文；其二，跳跃连接将编码器各层的局部细节直接传递至解码器对应层，补偿了上采样过程中的信息损失；其三，编码器-解码器的对称设计使解码器的特征分辨率与编码器对齐，避免了直接从最深层特征恢复全分辨率输出时的语义-分辨率失配问题。

\begin{figure}[htbp]
	\centering
	\includegraphics[width=\linewidth]{architecture.jpg}
	\caption{PointSS整体架构，括号内标注（点数，特征维度）表示各阶段张量形状}
	\label{fig:architecture}
\end{figure}

\section{输入处理与序列化}

原始点云输入包含三维坐标$(x, y, z)$及可选的RGB颜色信息。输入处理阶段包含数据增强和点云序列化两个步骤。数据增强在训练阶段施加随机翻转、随机缩放、随机旋转和点丢弃等操作以提升模型鲁棒性，推理阶段不使用增强。

序列化采用PTv3~\cite{ptv3}的空间填充曲线排列策略，将无序三维点云映射为一维有序序列。PTv3同时支持Z-order曲线与Hilbert曲线两种编码，Hilbert曲线具有更好的局部性保持特性，空间上相邻的点在序列中倾向于保持相邻关系，从而为后续的SSM序列建模提供合理的局部上下文。序列化仅在训练开始时执行一次，后续的下采样通过对序列化点云直接池化实现，避免了每个下采样阶段重复执行$O(N\log N)$序列化的开销。

\textbf{空间填充曲线编码。}以Z-order编码为例说明编码原理。在栅格深度$d$下，将每个点的坐标量化为$d$位整数，再将$x, y, z$三个维度的二进制位交织排列，得到$3d$位的编码值，排序即为序列化顺序。以$d=2$、坐标$(x, y, z)=(2, 2, 2)$为例：
\begin{align*}
x = 2 = 10_2, \quad y = 2 = 10_2, \quad z = 2 = 10_2
\end{align*}
将各维度二进制位按$z_1 y_1 x_1\, z_0 y_0 x_0$的顺序交织排列（下标表示位序，1为高位，0为低位），得到6位编码$111000_2 = 56$。对其邻居$(3, 2, 2)$，$x$的最低位由0变为1，交织后末位变化，编码为$111001_2 = 57$。两者按编码排序后在序列中紧邻，体现了空间邻近性。截断最低3位后两者同为$111_2 = 7$，即归入同一体素簇，这正是下文串行池化的原理基础。

\section{层次化编码器}

编码器由五个阶段构成，各阶段的通道维度依次为$(32, 64, 128, 256, 512)$，每个阶段的块数依次为$(2, 2, 2, 6, 2)$，多头注意力头数依次为$(2, 4, 8, 16, 32)$，点数在各下采样阶段依次减半。每个编码阶段（除第一阶段外）由一个下采样层和若干特征学习块串联而成。

在进入编码器前，输入特征首先经过嵌入层处理。嵌入层采用步长为1的稀疏三维卷积（SubMConv3d，核大小为$5\times5\times5$）将原始$C_{\text{in}}$维输入特征（坐标与颜色共6维）投影到第一阶段的通道维度32。稀疏卷积在体素化的离散坐标网格上操作，利用点云的稀疏性跳过空体素，兼顾了局部空间感知和计算效率。若启用GGAM模块，嵌入层在稀疏卷积之前还会注入几何增强特征，详见第四章。

\textbf{基于编码截断的下采样。}串行池化（SerializedPooling）的核心思想是对空间填充曲线编码进行低位截断，直接在编码空间中完成空间聚合，无需额外的三维近邻搜索。

给定点云在栅格深度$d$下的空间填充曲线编码$o_i$（$3d$位整数，包含$x, y, z$各$d$位的交织编码），其中最高位编码最粗粒度的空间划分，最低位编码最细粒度的空间划分。由于$x, y, z$三个维度的位按序交织排列，编码中最低的3位（即每个维度各贡献最低1位）共同描述了点在当前分辨率下所在的最细粒度体素格。设$d_{\text{pool}}$为每次池化截断的精度层数（本文取$d_{\text{pool}}=1$，对应下采样步长为2），将编码右移$3d_{\text{pool}}$位，等价于同时去掉每个维度最低$d_{\text{pool}}$位精度，使原本仅在这些低位上有差异（即在空间上相邻）的点获得相同的截断编码：
\begin{equation}
    o_i' = o_i \gg 3 d_{\text{pool}}
\end{equation}
截断后，共享同一编码$o_i'$的点必然落在空间上相邻的体素内，无需任何三维近邻搜索即可将其归入同一簇$k$。在几何意义上，这等价于将$2^{d_{\text{pool}}} \times 2^{d_{\text{pool}}} \times 2^{d_{\text{pool}}} = 2\times2\times2$个相邻体素合并为一个粗粒度体素。每个簇的代表特征和坐标分别为：
\begin{equation}
    \mathbf{f}_k = \underset{i:\,o_i'=k}{\text{MaxPool}}\!\left(\mathbf{W}_{\text{proj}}\mathbf{f}_i\right), \quad
    \mathbf{p}_k = \underset{i:\,o_i'=k}{\text{MeanPool}}(\mathbf{p}_i)
\end{equation}
其中$\mathbf{W}_{\text{proj}}$为可学习的线性投影矩阵，在聚合前将特征从$C_s$维变换到$C_{s+1}$维。特征采用最大池化而坐标采用均值池化，兼顾了特征的判别性与位置的几何一致性。串行池化的计算复杂度为$O(N)$，且全程复用已有的序列化编码，无需额外空间查询。下采样时还记录每个细粒度点所属的簇索引$\sigma(i)=k$，供解码阶段逆映射使用。

\textbf{特征学习块与三重残差结构。}每个编码阶段内，特征学习块负责在当前分辨率下对点云特征进行深度建模。每个块包含三条独立的残差路径（如\cref{fig:coder}所示），依次作用于输入特征$\mathbf{f}$：
\begin{equation}
    \mathbf{h}_0 = \mathbf{f} + \text{CPE}(\mathbf{f})
    \label{eq:res_cpe}
\end{equation}
\begin{equation}
    \mathbf{h}_1 = \mathbf{h}_0 + \text{DropPath}\!\left(\mathcal{F}_{\text{SSM}}\!\left(\text{LN}(\mathbf{h}_0)\right)\right)
    \label{eq:res_ssm}
\end{equation}
\begin{equation}
    \mathbf{h}_2 = \mathbf{h}_1 + \text{DropPath}\!\left(\text{MLP}\!\left(\text{LN}(\mathbf{h}_1)\right)\right)
    \label{eq:res_mlp}
\end{equation}
其中$\text{LN}(\cdot)$为层归一化，$\mathcal{F}_{\text{SSM}}(\cdot)$为ASD-SSM（和/或Chebyshev多尺度SSM）模块，$\text{DropPath}(\cdot)$为随机深度正则化操作。

第一条残差路径（\cref{eq:res_cpe}）引入条件位置编码（Conditional Positional Encoding, CPE）。CPE由稀疏三维卷积（SubMConv3d，核大小$3\times3\times3$）、线性层和层归一化串联构成，根据点在稀疏体素邻域内的实际分布动态生成位置偏置，使模型对不同密度和布局的点云具有更强的适应性。后两条残差路径分别对应SSM建模和MLP特征变换，均采用前置归一化（Pre-Norm）结构——先归一化再计算，有助于在深层网络中稳定梯度流。

DropPath（随机深度）以概率$p_l$在训练时将第$l$块的残差分支输出置零，其概率按层线性递增：
\begin{equation}
    p_l = \frac{l}{L} \cdot p_{\max}
\end{equation}
其中$L$为编码器总块数，$p_{\max}$为设定的最大丢弃率。这一渐进式正则化策略在保留浅层特征的同时对深层特征施加更强的随机扰动，有效缓解了深层网络的过拟合。

\textbf{顺序提示与多序列化轮换。}为使SSM能够感知当前处理的序列化方向，每个特征学习块在调用SSM前会将$P=6$个可学习的顺序提示令牌（Order Prompt Token）分别前置和后置于点的特征序列，SSM处理完毕后再将其剥离。PointSS支持四种序列化顺序（Z序、转置Z序、Hilbert序、转置Hilbert序），每种顺序维护一组独立的提示令牌，通过线性层投影至当前层的通道维度。在编码器各阶段内，相邻特征学习块循环使用不同的序列化顺序，使网络在不同扫描方向下均能建模序列依赖，有效缓解了单一序列化方向下空间覆盖不完整的问题。

\begin{figure}[htbp]
	\centering
	\includegraphics[width=\linewidth]{coder.JPG}
	\caption{编码器基本单元结构。池化层（SerializedPooling）通过编码截断进行下采样，多个特征学习块（Block）以三重残差结构进行深层特征提取}
	\label{fig:coder}
\end{figure}

\section{解码器与跳跃连接}

解码器与编码器对称，由四个上采样阶段构成，通道维度依次为$(64, 64, 128, 256)$，每个阶段的块数为$(2, 2, 2, 2)$。每个解码阶段由上采样层（串行反池化）和若干特征学习块串联构成，特征学习块与编码器共享相同的三重残差结构（\crefrange{eq:res_cpe}{eq:res_mlp}）。

\textbf{基于簇索引的上采样。}解码器的上采样操作（SerializedUnpooling）与编码器下采样形成精确的逆对应：编码器在SerializedPooling时为每个细粒度点记录了其所属的粗粒度簇索引$\sigma(i)$，上采样阶段直接利用该索引将粗粒度特征传播回各细粒度点，无需重新执行三维空间近邻搜索。

设$\mathbf{f}_i^{\text{enc}}$为编码器在该分辨率下存储的跳跃连接特征，$\mathbf{f}_k^{\text{dec}}$为解码器当前层的粗粒度特征，细粒度点$i$的上采样输出为：
\begin{equation}
    \mathbf{f}_i^{\text{up}} = \mathbf{W}_{\text{skip}}\,\mathbf{f}_i^{\text{enc}} + \mathbf{W}_{\text{dec}}\,\mathbf{f}_{\sigma(i)}^{\text{dec}}
    \label{eq:unpooling}
\end{equation}
其中$\mathbf{W}_{\text{skip}} \in \mathbb{R}^{C_{\text{out}} \times C_{\text{skip}}}$和$\mathbf{W}_{\text{dec}} \in \mathbb{R}^{C_{\text{out}} \times C_{\text{dec}}}$分别为跳跃连接路径和解码路径的线性投影矩阵，两路投影后特征以逐元素加法融合。加法融合（而非拼接）的优势在于：在不增加后续层输入维度的前提下实现了两路信息的叠加，避免了拼接方案中因通道数翻倍带来的参数和计算量增加。

\textbf{跳跃连接的信息互补性。}\cref{eq:unpooling}中的$\sigma(i)$由下采样阶段直接给出，因而粗粒度特征的传播本质上是零额外开销的索引查找。跳跃连接路径$\mathbf{W}_{\text{skip}}\mathbf{f}_i^{\text{enc}}$将编码器在该分辨率下建模的局部几何细节直接引入解码器，解码路径$\mathbf{W}_{\text{dec}}\mathbf{f}_{\sigma(i)}^{\text{dec}}$则携带了更深层编码器建模的场景级语义上下文。二者在各分辨率层级上逐步融合，使解码器在恢复空间分辨率的同时保留精细的边界和局部形状信息，有助于精细的逐点语义分割。

\section{任务预测头}

PointSS的预测头根据下游任务进行差异化设计，以适配语义分割和分类两类主要任务。

\textbf{语义分割头。}语义分割任务要求对每个点输出类别预测。解码器最终输出的全分辨率特征（维度64）经一个线性分类层直接映射为$K$维类别预测向量（$K$为类别数），训练时以交叉熵损失监督。交叉熵损失将softmax归一化后的预测概率与真实标签对比，对错误分类的点施加较大惩罚，适合类别互斥的语义分割场景。对于S3DIS数据集，$K=13$；对于ScanNet数据集，$K=20$。

\textbf{分类头。}点云分类任务要求对整体点云输出单一类别标签。编码器最深层的点云特征（通道512，点数约为原始点数的$1/32$）经全局平均池化压缩为固定维度的场景级表示，再经两层MLP（含Dropout和BatchNorm）映射至$K$维类别预测向量（ModelNet40中$K=40$），训练时同样采用交叉熵损失。

\section{本章小结}

本章系统阐述了PointSS的整体架构设计。PointSS基于U型编码器-解码器框架，以串行池化下采样和基于簇索引的逆映射上采样构建层次化的特征提取流程，通道维度从浅层的32逐步扩展至深层的512。序列化预处理复用PTv3的Hilbert曲线策略，且仅执行一次以规避重复序列化的计算开销。下采样通过对空间填充曲线编码进行低位截断实现$O(N)$的空间聚合，并记录簇索引供解码器逆映射使用，无需重新执行空间搜索。每个特征学习块包含CPE、SSM和MLP三条独立的残差路径，配合前置归一化和DropPath随机深度正则化，以及多序列化顺序轮换机制，保证了深层网络的训练稳定性与多方向的序列依赖建模。解码器的跳跃连接通过加法融合的方式，在各分辨率层级上将编码器的局部细节与解码器的语义上下文叠加，保留了精细的边界和形状信息。针对不同下游任务，分别设计了基于逐点线性分类的分割头和基于全局池化MLP的分类头。在此整体框架之上，后续各章将分别介绍嵌入其中的三个创新模块：GGAM、ASD-SSM和Chebyshev多尺度SSM。

\chapter{基于全局几何感知的点云特征增强}
基于Mamba的点云分析方法面临局部几何结构建模困难的问题，根本原因在于序列化状态传递机制与无序三维结构之间的冲突。Mamba通过以下离散化状态空间方程进行序列建模：
\begin{equation}
	x_k = \bar{A}x_{k-1} + \bar{B}u_k, \quad y_k = \bar{C}x_k + \bar{D}u_k
\end{equation}
其中$k$表示序列化后的位置索引，$x_k$为第$k$个位置的状态向量，$u_k$为输入特征向量，$y_k$为输出特征向量，$\bar{A}$、$\bar{B}$、$\bar{C}$、$\bar{D}$分别为状态转移矩阵、输入投影矩阵、输出投影矩阵和跳跃连接矩阵。

在该递归机制中，每个状态$x_k$依赖于前一状态$x_{k-1}$和当前输入$u_k$，从而建立序列依赖关系。如上文所述，序列化不可避免地破坏了空间邻近性——空间上相邻的点在序列中有可能相隔很远。这种邻近性的丧失使得递归状态更新无法有效聚合空间邻居的信息，限制了模型捕获局部几何结构（如法向量和曲率）的能力。

针对上述问题，我们提出了全局几何感知机制（Global Geometry-Aware Mechanism, GGAM），通过显式注入几何先验来补偿序列化造成的几何信息丢失。如\cref{fig:ggam}所示，GGAM由两部分组成：（1）图构建与几何特征提取；（2）几何特征增强与融合。第一部分从序列化的点序列中构建图结构并提取每个点的几何特征向量。为捕获互补的空间结构信息，该过程在Z-order和Hilbert两种序列化流上并行执行，产生两组独立的几何特征。第二部分通过交叉注意力机制和自适应门控策略对双流特征进行增强与融合。后续将围绕方法设计和相关消融实验进行介绍。

\begin{figure}[htbp]
	\centering
	\includegraphics[width=\linewidth]{ggam.jpg}
	\caption{全局几何感知机制（GGAM）的整体结构}
	\label{fig:ggam}
\end{figure}

\section{方法设计}
\textbf{图构建与几何特征提取。}传统的几何关系计算方法（如曲率、法向量）依赖$K$近邻（KNN）搜索来构建邻域。然而，KNN搜索的计算复杂度较高，难以适应大规模点云场景。受Point Transformer V3\cite{ptv3}的启发——该工作表明基于序列化的窗口邻域可以替代KNN，在不损失性能的前提下将复杂度从$\mathcal{O}(N\log N)$降至$\mathcal{O}(N)$——我们采用该策略来实现高效的几何特征提取。

给定点云$\{p_1, p_2, \ldots, p_N\}$，序列化方法按照空间填充曲线对其进行排序，得到序列$Q = (q_1, q_2, \ldots, q_N)$，其中$q_i$表示排序后第$i$个位置上的点。随后将$Q$划分为不重叠的窗口：
\begin{equation}
	W_m = (q_{(m-1) \cdot 64 + 1}, q_{(m-1) \cdot 64 + 2}, \ldots, q_{m \cdot 64})
\end{equation}
其中$m = 1, 2, \ldots, \lfloor N/64 \rfloor$为窗口索引，窗口大小固定为64，该尺度能为基于PCA的法向量和曲率估计提供充足的邻域。若最后一个窗口不足64个点，则按照PTv3\cite{ptv3}的策略进行填充。

对于每个点$q_i$，其邻域$\mathcal{N}(q_i)$由同一窗口内的所有其他点组成，构成全连接的局部图：
\begin{equation}
	\mathcal{N}(q_i) = \left\{q_j : q_j \in W_{\lceil i/64 \rceil},\ j \neq i\right\}
\end{equation}

对于每对点$(q_i, q_j)$（$q_j \in \mathcal{N}(q_i)$），提取8维几何特征向量：
\begin{equation}
	\mathbf{f}_{ij} = [\Delta q_{ij},\ \mathbf{n}_i \cdot \mathbf{n}_j,\ \cos\theta_i,\ \cos\theta_j,\ \kappa_i,\ \kappa_j] \in \mathbb{R}^8
\end{equation}
其中$\Delta q_{ij} = q_j - q_i \in \mathbb{R}^3$为相对位置向量；$\mathbf{n}_i, \mathbf{n}_j \in \mathbb{R}^3$为$q_i$和$q_j$处的法向量，通过对窗口内点的三维坐标进行PCA估计得到。$\mathbf{n}_i \cdot \mathbf{n}_j \in \mathbb{R}$为法向量的点积，用于度量两端点之间的法向一致性——在光滑表面上趋近于1，在几何边界处急剧下降。$\cos\theta_i$和$\cos\theta_j$分别为相对位置向量$\Delta q_{ij}$与法向量$\mathbf{n}_i$和$\mathbf{n}_j$之间夹角的余弦值：
\begin{equation}
	\cos\theta_i = \frac{\mathbf{n}_i \cdot \Delta q_{ij}}{\|\Delta q_{ij}\|}, \quad \cos\theta_j = \frac{\mathbf{n}_j \cdot \Delta q_{ij}}{\|\Delta q_{ij}\|}
\end{equation}
$\cos\theta_i$度量边方向$\Delta q_{ij}$偏离$q_i$处切平面的程度：在光滑区域接近零，在几何边缘处显著增大，$\cos\theta_j$在$q_j$处类似定义。$\kappa_i, \kappa_j \in \mathbb{R}$为两端点处的表面曲率。

上述基于法向量的特征分量——$\mathbf{n}_i \cdot \mathbf{n}_j$、$\cos\theta_i$、$\cos\theta_j$、$\kappa_i$和$\kappa_j$——均为旋转不变的几何量\cite{clusterNet, riconv}：对点云施加全局旋转不会改变这些数值。相比之下，原始法向量会随点云旋转而变化，不具备这种不变性。旋转不变性有助于增强模型对旋转变化的鲁棒性\cite{clusterNet, riconv}。

每条边特征经过多层感知机（MLP）处理后，对每个点$q_i$通过其邻域$\mathcal{N}(q_i)$上的均值池化进行聚合：
\begin{equation}
	\mathbf{H}_i = \frac{1}{|\mathcal{N}(q_i)|} \sum_{q_j \in \mathcal{N}(q_i)} \text{MLP}(\mathbf{f}_{ij})
\end{equation}
所得特征$\mathbf{H}_i \in \mathbb{R}^C$编码了$q_i$的局部几何上下文信息。由于Z-order序列化偏向于局部聚集性，而Hilbert序列化更好地保持空间连续性，图构建在两种序列化上并行执行，为每个点生成两组互补的几何特征向量$\mathbf{H}_z$和$\mathbf{H}_h$，并传递给GGAM的第二部分。

\textbf{几何特征增强与融合。}几何特征增强与融合模块对$\mathbf{H}_z$和$\mathbf{H}_h$依次进行三个步骤的处理：流内自注意力、跨流交叉注意力和自适应门控融合。

首先在每个流内部施加自注意力，其中$\text{Attention}(\mathbf{Q}, \mathbf{K}, \mathbf{V}) = \text{Softmax}(\mathbf{Q}\mathbf{K}^\top / \sqrt{d})\mathbf{V}$：
\begin{equation}
	\begin{aligned}
		\mathbf{A}_z = \text{Attention}(\mathbf{H}_z, \mathbf{H}_z, \mathbf{H}_z), \\
		\mathbf{A}_h = \text{Attention}(\mathbf{H}_h, \mathbf{H}_h, \mathbf{H}_h)
	\end{aligned}
\end{equation}
并通过残差连接得到：
\begin{equation}
	\mathbf{H}'_z = \mathbf{H}_z + \mathbf{A}_z, \quad \mathbf{H}'_h = \mathbf{H}_h + \mathbf{A}_h
\end{equation}
然后施加交叉注意力，每个流保留自身的查询向量，而键和值则来自另一个流：
\begin{equation}
	\begin{aligned}
		\mathbf{A}'_z = \text{Attention}(\mathbf{H}'_z, \mathbf{H}'_h, \mathbf{H}'_h), \\
		\mathbf{A}'_h = \text{Attention}(\mathbf{H}'_h, \mathbf{H}'_z, \mathbf{H}'_z)
	\end{aligned}
\end{equation}
通过残差连接得到增强后的特征：
\begin{equation}
	\mathbf{H}''_z = \mathbf{H}'_z + \mathbf{A}'_z, \quad \mathbf{H}''_h = \mathbf{H}'_h + \mathbf{A}'_h
\end{equation}
为自适应地融合两组增强特征，自适应门控单元为每个流生成一个标量权重：
\begin{equation}
	[\alpha_z, \alpha_h] = \text{Softmax}(\text{MLP}([\mathbf{H}''_z \| \mathbf{H}''_h]))
	\label{eq:adaptive_gate}
\end{equation}
其中$\|$为拼接操作。融合特征通过加权求和得到：
\begin{equation}
	\mathbf{H}_{\text{fused}} = \alpha_z \mathbf{H}''_z + \alpha_h \mathbf{H}''_h
	\label{eq:gated_fusion}
\end{equation}
其中$\alpha_z, \alpha_h \in \mathbb{R}$为标量权重且满足$\alpha_z + \alpha_h = 1$，在所有通道上统一广播。

由于两个流建立在不同序列化产生的窗口划分之上，其融合表示可能包含跨流不一致性。为缓解该问题，几何一致性约束模块从两个流中推导出一个可学习的标量门控$\gamma \in [0, 1]$，用于自适应地调制融合表示：
\begin{equation}
	\gamma = \text{Sigmoid}(\text{MLP}([\mathbf{H}''_z \| \mathbf{H}''_h]))
	\label{eq:consistency_factor}
\end{equation}
\begin{equation}
	\mathbf{H}_{\text{final}} = \mathbf{H}_{\text{fused}} \odot \gamma
	\label{eq:final_feature}
\end{equation}
$\mathbf{H}_{\text{final}}$作为GGAM模块的输出，传递给后续的编码器和解码器，提供几何感知的点云表示。

\section{实验设置与分析}
本节通过系统性的消融实验验证GGAM在缓解点云序列化所引入的空间邻近性丢失问题上的有效性。消融实验在S3DIS数据集上开展，通过三组实验分别验证GGAM的整体补偿能力（\cref{tab:ggam_fragmentation}）、逐类别补偿效果（\cref{tab:ggam_category}）以及各组件的独立贡献（\cref{tab:ggam_components}）。基线模型采用普通单尺度双向Mamba架构，使用共享的单一状态转移参数，保持与完整PointSS相同的网络深度和通道配置。详细的训练配置与超参数设置将在第六章实验部分统一介绍。

为定量衡量序列化后每个点的空间邻近性丢失程度，我们以每个点为中心搜索其128个空间近邻点，然后计算这些近邻点在序列化后与中心点距离超过300个位置的近邻点所占比例。选择300作为距离阈值是因为其远大于窗口大小（64），能够提供更宽泛的采样范围以获得更稳定的邻近性丢失估计。我们将这一逐点度量称为\textit{邻近性丢失比率}，该比率越高，表明该点局部邻域的空间邻近性丢失越严重。我们据此将点按照空间邻近性丢失程度进行分组，并比较基线模型与GGAM的分割性能，结果如\cref{tab:ggam_fragmentation}所示。

\begin{table}[htbp!]
	\centering
	\caption{按邻近性丢失程度分层的GGAM效果分析}
	\label{tab:ggam_fragmentation}
	\begin{tabular}{@{}llllll@{}}
		\toprule
		邻近性丢失比率 & 点数量 & 点占比 & 基线IoU (\%) & GGAM IoU (\%) & 绝对提升 (\%) \\
		\midrule
		0-10\%（低） & 27.0万 & 67\% & 74.0 & 74.8 & 0.8 \\
		10-20\%（轻度） & 10.1万 & 25\% & 63.0 & 67.1 & 4.1 \\
		20-30\%（中度） & 3.2万 & 8\% & 62.0 & 70.0 & 8.0 \\
		\midrule
		总计 & 40.3万 & 100\% & 70.3 & 72.5 & 2.2 \\
		\bottomrule
	\end{tabular}
	\par\vspace{3pt}
	\footnotesize 注：邻近性丢失比率 = 该点在空间内的128个近邻点在序列中与中心点距离超过300个位置的点数 / 临近点总数（128）。
\end{table}

实验结果表明，在Hilbert曲线序列化下，约33\%的点存在不同程度的空间邻近性丢失。随着邻近性丢失比率从0--10\%增加到20--30\%，基线模型性能急剧下降，证实了序列化对局部几何表征的影响。GGAM有效缓解了这一性能退化，尤其在高丢失比率组（20--30\%）中获得了显著的性能提升，远超低丢失比率组的改善幅度。

逐类别分析（\cref{tab:ggam_category}）进一步证实，邻近性丢失比率较高的类别整体上从GGAM中获益更多。大平面结构（如地板、墙面）受序列化引发的空间邻近性丢失影响较小，性能提升有限。随着邻近性丢失比率的增加，GGAM的补偿效果更为显著：具有复杂边界或不规则结构的类别（如窗户、门、板、杂物）表现出更大的性能提升。

\begin{table}[htbp!]
	\centering
	\caption{S3DIS数据集上GGAM的逐类别性能分析}
	\label{tab:ggam_category}
	\begin{tabular}{@{}llllll@{}}
		\toprule
		类别 & 平均邻近性丢失比率 & 基线IoU (\%) & GGAM IoU (\%) & 绝对提升 (\%) \\
		\midrule
		地板   & 5\%    & 97.3  & 98.5  & +1.2 \\
		墙     & 6\%    & 84.4  & 85.8  & +1.4 \\
		天花板 & 8\%    & 89.9  & 92.3  & +2.4 \\
		桌子   & 11\%   & 80.8  & 82.5  & +1.7 \\
		书架   & 13\%   & 78.5  & 80.0  & +1.5 \\
		椅子   & 15\%   & 91.1  & 93.6  & +2.5 \\
		沙发   & 15\%   & 81.5  & 84.3  & +2.8 \\
		板     & 15\%   & 84.8  & 87.9  & +3.1 \\
		门     & 16\%   & 77.6  & 80.6  & +3.0 \\
		柱子   & 18\%   & 33.6  & 35.4  & +1.8 \\
		窗户   & 18\%   & 58.9  & 62.4  & +3.5 \\
		杂物   & 22\%   & 56.1  & 59.4  & +3.3 \\
		横梁   & --     & 0.0   & 0.0   & 0.0 \\
		\midrule
		平均   & 13.5\% & 70.3  & 72.5  & +2.2 \\
		\bottomrule
	\end{tabular}
	\par\vspace{3pt}
	\footnotesize 注：平均邻近性丢失比率为特定语义类别中所有点的邻近性丢失比率均值。横梁类别由于在S3DIS Area 5中实例极为稀少，两种模型的召回率均接近零，因此未纳入分析。
\end{table}

我们进一步对GGAM的各组件进行消融实验（\cref{tab:ggam_components}），以验证其各自的贡献。单独引入双序列化会导致轻微的性能下降，这是因为不同序列化方式引入了互补但可能冲突的邻域偏置，在缺乏有效融合策略的情况下难以协调。加入跨序列注意力后性能获得明显提升，表明跨序列交互在对齐不同序列化模式下的互补几何信息方面发挥了重要作用。引入门控融合机制后，通过自适应平衡两个序列化分支的贡献，进一步优化了融合后的表征。加入几何一致性门控后又获得了额外提升，该门控从两个流中推导标量门控信号来调制融合输出，缓解了跨流不一致性问题。

\begin{table}[htbp!]
	\centering
	\caption{S3DIS数据集上GGAM的组件消融实验}
	\label{tab:ggam_components}
	\begin{tabular}{@{}ccccc@{}}
		\toprule
		DS & CSA & GFM & $\mathcal{L}_{geo}$ & mIoU (\%) \\
		\midrule
		$\times$ & $\times$ & $\times$ & $\times$ & 71.1 \\
		$\checkmark$ & $\times$ & $\times$ & $\times$ & 70.6 \\
		$\checkmark$ & $\checkmark$ & $\times$ & $\times$ & 71.9 \\
		$\checkmark$ & $\checkmark$ & $\checkmark$ & $\times$ & 72.2 \\
		$\checkmark$ & $\checkmark$ & $\checkmark$ & $\checkmark$ & 72.5 \\
		\bottomrule
	\end{tabular}
	\par\vspace{3pt}
	\footnotesize 注：DS：双序列化；CSA：跨序列注意力；GFM：门控融合机制；$\mathcal{L}_{geo}$：几何一致性约束（门控模块）。所有配置均保留几何特征提取，因其为GGAM的前提条件。此处基线（71.1\%）与\cref{tab:progressive_improvement}中的纯Mamba基线（70.3\%）不同，因为所有配置均保留了几何特征提取。
\end{table}

完整组合表明，这些组件作为一个整体协同运作，共同缓解了序列化引发的空间邻近性丢失问题。


\section{本章小结}

本章针对序列化状态传递机制与无序三维结构之间的冲突所导致的几何信息丢失问题，提出了全局几何感知机制（GGAM）。本章首先分析了该问题的根本原因：序列化不可避免地破坏了空间邻近性，空间上相邻的点在序列中可能相隔很远，使得递归状态更新无法有效聚合空间邻居的信息，限制了模型捕获局部几何结构的能力。针对该问题，GGAM通过窗口化邻域高效提取相对位置、法向量和曲率等几何先验，并在Z-order和Hilbert两种序列化流上并行执行图构建与几何特征提取，产生两组互补的几何特征。在特征增强与融合阶段，GGAM依次通过流内自注意力、跨流交叉注意力和自适应门控融合对双流特征进行增强与融合，并引入几何一致性约束缓解跨流不一致性。消融实验表明，GGAM能够有效缓解序列化引发的空间邻近性丢失，且邻近性丢失程度越高的点和几何结构越复杂的类别，从GGAM中获得的性能补偿越显著。组件消融进一步验证了双序列化、跨序列注意力、门控融合和几何一致性约束四个组件协同运作的必要性。GGAM作为PointSS的基础模块，其输出的几何增强特征贯穿整个编码-解码过程，与后续章节提出的ASD-SSM和Chebyshev多尺度SSM共同构成完整的技术体系。


\chapter{多尺度点云状态空间建模}

点云序列建模的核心挑战在于如何在单一模块中同时捕获不同尺度的几何信息。本章提出两种互补的多尺度建模方法：ASD-SSM在序列化的点云上通过自适应尺度解耦设计显式构建多尺度特征；Chebyshev多尺度SSM则借助图拉普拉斯的多项式展开，将点云特征分解为不同跳数邻域的结构分量，并由独立的Mamba模块分别建模。两者分别从序列建模和图结构建模的角度形成互补，共同构成本文的多尺度特征学习框架。

\section{基于序列化的多尺度建模：ASD-SSM}

本节提出自适应尺度解耦状态空间模型（Adaptive Scale-Decoupled State Space Model, ASD-SSM）。ASD-SSM通过引入动态参数生成机制来解决状态转移参数输入无关的局限性，使$\overline{A}$能够根据每个Patch的局部几何内容进行自适应调整。这种尺度解耦设计使细尺度能够捕获几何细节，而粗尺度建模语义上下文，实现有效的层次化表示。

在后续描述中，$s \in \{1, 2, \ldots, S\}$表示尺度索引，其中$s=1$表示最细尺度，$s=S$表示最粗尺度。在尺度$s$下，点云被划分为$M_s$个不重叠的Patch，$m \in \{1, \ldots, M_s\}$表示Patch索引。$N$和$C$分别表示点的数量和特征维度。

\begin{figure}[htbp]
	\centering
	\includegraphics[width=\linewidth]{BSS.JPG}
	\caption{ASD-SSM具体流程}
	\label{fig:bss}
\end{figure}


如\cref{fig:bss}所示，ASD-SSM的处理流程包含三个阶段：窗口划分、顺序提示和多尺度特征建模。

\subsection{窗口划分}

序列化后的点云被均匀划分为不重叠的Patch。若最后一个Patch的点数不足指定大小，则按照PTv3~\cite{ptv3}的策略进行填充。在尺度$s$下，Patch大小为$P \times F_s$（最细尺度$F_1 = 1$）；对于三尺度配置（$S=3$），扩展因子分别为$F_2=2$和$F_3=4$。

\subsection{顺序提示}

参照PCM\cite{pcm}的设计，顺序提示在最细尺度的每个Patch首尾插入可学习的提示向量，告知Mamba当前序列的排列规则；在较粗尺度下，每个Patch内的点被随机打乱且不插入提示。在粗尺度下，建模目标从局部几何结构转向全局语义上下文。若在粗尺度保留序列化顺序，模型可能利用局部序列模式作为捷径，导致粗尺度特征无意中依赖局部排列顺序而非真正聚合全局语义信息。打乱操作消除了这一捷径，迫使模型仅从特征内容本身构建全局表示。


\subsection{基于尺度自适应参数生成的多尺度特征建模}

虽然现有基于SSM的方法通过选择性机制实现了输入相关的$\overline{B}$和$\overline{C}$，但状态转移矩阵$A$在训练后保持固定，所有Patch共享同一组$A$参数。步长参数$\Delta$虽然是输入相关的，通过离散化$\overline{A} = \exp(\Delta A)$使$\overline{A}$逐点变化，但无法根据不同Patch的几何内容和所处尺度自适应地调整状态转移特性。$\overline{A}$的对角元素控制状态衰减速率：接近1时实现慢衰减（长程记忆），接近0时实现快衰减（快速响应局部变化）。因此，粗尺度特征需要接近1的$\overline{A}$值以捕获长程依赖，而细尺度特征需要较小的值以快速响应局部几何细节。

\textbf{尺度感知参数生成器：}对于给定尺度$s$，点云特征被划分为Patch，参数生成过程对每个Patch独立进行，根据Patch的内容特征和尺度信息生成专属的状态转移参数$\overline{A}_s$。

具体而言，对于尺度$s$的第$m$个Patch，我们对Patch内所有点进行全局平均池化得到该Patch的全局特征表示$\mathbf{f}_{\text{global}}^{(s,m)} \in \mathbb{R}^{C}$，然后通过$\text{MLP}_{\text{global}}$投影至$\mathbb{R}^{C/4}$：
\begin{equation}
	\tilde{\mathbf{f}}_{\text{global}}^{(s,m)} = \text{MLP}_{\text{global}}(\mathbf{f}_{\text{global}}^{(s,m)}) \in \mathbb{R}^{C/4}
\end{equation}
同时为每个尺度$s$维护一个可训练的尺度嵌入$\mathbf{e}_s \in \mathbb{R}^{C/4}$，作为尺度标识符。

\textbf{尺度约束机制：}为实现差异化的状态衰减特性，通过尺度约束因子$\alpha_s$显式控制$\overline{A}_s$的值域。在离散状态空间模型框架中，状态更新遵循$\mathbf{h}_k = \overline{A}_s\mathbf{h}_{k-1} + \overline{B}_s\mathbf{u}_k$。

为直接约束$\overline{A}_s$的值域，首先通过sigmoid函数将参数生成器的输出映射到$[0,1]$区间：
\begin{equation}
	\hat{A}_{s,m}^{\text{raw}} = \sigma(\lambda \cdot \text{MLP}([\tilde{\mathbf{f}}_{\text{global}}^{(s,m)} \| \mathbf{e}_s]))
	\label{eq:param_norm}
\end{equation}
其中$\sigma(\cdot)$为sigmoid激活函数，$\|$表示拼接，$\lambda$为温度参数（设为0.1），用于控制MLP输出经过sigmoid激活前的灵敏度。然后，尺度约束因子$\alpha_s$作为上界直接控制$\overline{A}_s$的值域：
\begin{equation}
	\overline{A}_{s,m} = \alpha_s \cdot \hat{A}_{s,m}^{\text{raw}}
	\label{eq:param_modulation}
\end{equation}

通过设置$\alpha_s$从细尺度到粗尺度单调递增（$\alpha_1 = 0.3$到$\alpha_S = 0.9$），我们施加显式约束：细尺度$\overline{A}_1 \in [0, 0.3]$实现快速响应局部变化，粗尺度$\overline{A}_S \in [0, 0.9]$保持长程记忆以捕获全局上下文。需要强调的是，在每个尺度内，$\overline{A}_s$根据每个Patch的局部内容独立生成，而非在整个尺度上共享。对于选择性参数$\overline{B}_s$和$\overline{C}_s$，沿用Mamba\cite{Mamba}的标准选择性机制，通过线性投影从输入内容动态生成。

\textbf{状态空间更新：}在每个Patch内，状态更新遵循标准的离散化SSM递推形式\cite{Mamba}：
\begin{equation}
	\begin{aligned}
		\mathbf{h}_k^{(s,m)} &= \overline{A}_{s,m} \mathbf{h}_{k-1}^{(s,m)} + \overline{B}_s \mathbf{u}_k^{(s,m)} \\
		\mathbf{y}_k^{(s,m)} &= \overline{C}_s \mathbf{h}_k^{(s,m)} + \overline{D} \mathbf{u}_k^{(s,m)}
	\end{aligned}
	\label{eq:asd_ssm_update}
\end{equation}
其中$\mathbf{h}_k^{(s,m)}$是尺度$s$、Patch $m$中第$k$个位置的隐状态向量，$\mathbf{h}_{k-1}^{(s,m)}$为前一位置的隐状态，$\mathbf{u}_k^{(s,m)}$是第$k$个点的输入特征向量，$\overline{D}$为直通连接参数。为保持跨Patch边界的序列连续性，每个Patch的最终隐状态被传递为下一个Patch的初始状态，而非重置为零向量，使得模型将整个序列化的点云序列视为连续流进行处理。

\textbf{双向特征学习与融合：}每个尺度内采用双向特征学习策略。我们为每个尺度$s$构建前向SSM和后向SSM，分别沿正向和反向处理序列化的点云序列，使每个点能够聚合来自两个方向的上下文信息。前向和后向SSM共享同一组尺度特定的状态转移参数$\overline{A}_s$，双向特征通过残差连接融合生成该尺度的输出。

通过上述机制，ASD-SSM实现了尺度解耦的状态空间建模：细尺度保持较小的$\overline{A}$值以实现快速状态衰减，捕获局部几何细节；粗尺度保持较大的$\overline{A}$值以实现缓慢衰减，建模长程全局上下文。所有尺度的特征提取完成后，将各尺度特征$\text{Points}_{1:S}$按通道维度逐点拼接，再通过线性投影和MLP映射回原始维度$C$，得到多尺度融合的点云表示$\text{Points}' \in \mathbb{R}^{N \times C}$。

\subsection{跨Patch状态传递机制}

为保证序列建模的连续性，ASD-SSM采用跨Patch状态传递机制。传统的窗口化处理方法通常在每个窗口边界处重置隐状态为零向量，这会导致窗口间的上下文信息丢失。ASD-SSM通过以下机制保持状态连续性：

**状态传递策略**：对于尺度$s$中的第$m$个Patch，其初始隐状态$\mathbf{h}_0^{(s,m)}$设置为前一个Patch的最终隐状态$\mathbf{h}_{P \times F_s}^{(s,m-1)}$。这种设计使得整个序列化点云被视为一个连续的数据流，而非独立的窗口片段。

**边界效应缓解**：由于不同Patch可能具有不同的几何特性，直接传递隐状态可能引入不一致性。为缓解这一问题，在状态传递时引入可学习的边界调制因子$\beta_s \in [0,1]$：
\begin{equation}
\mathbf{h}_0^{(s,m)} = \beta_s \cdot \mathbf{h}_{P \times F_s}^{(s,m-1)} + (1-\beta_s) \cdot \mathbf{0}
\end{equation}
其中$\beta_s$通过训练自适应学习，平衡状态连续性与Patch独立性之间的权衡。

**计算复杂度分析**：跨Patch状态传递的额外计算开销为$O(C)$，相比Patch内的状态更新复杂度$O(P \times F_s \times C)$可以忽略不计。这种设计在几乎不增加计算成本的前提下，显著提升了序列建模的连续性。

\subsection{多尺度特征融合策略}

ASD-SSM生成的多尺度特征需要有效融合以形成统一的点云表示。本文采用层次化融合策略：

**特征对齐**：由于不同尺度的Patch大小不同，需要将粗尺度特征上采样至细尺度的分辨率。对于尺度$s$的特征$\mathbf{F}_s \in \mathbb{R}^{N_s \times C}$，通过最近邻插值将其扩展至$\mathbf{F}_s' \in \mathbb{R}^{N \times C}$，其中$N$为原始点云大小。

**注意力加权融合**：不同尺度特征的重要性因点的局部几何特性而异。引入自注意力机制计算各尺度特征的权重：
\begin{equation}
\mathbf{w}_i = \text{Softmax}(\text{MLP}([\mathbf{F}_1'[i] \| \mathbf{F}_2'[i] \| \cdots \| \mathbf{F}_S'[i]]))
\end{equation}
最终融合特征为：
\begin{equation}
\mathbf{F}_{\text{fused}}[i] = \sum_{s=1}^{S} \mathbf{w}_i[s] \cdot \mathbf{F}_s'[i]
\end{equation}

这种自适应融合策略使模型能够根据每个点的局部特性动态调整不同尺度特征的贡献，提升了多尺度表示的表达能力。

\subsection{消融实验}

通过系统消融实验，分析自适应尺度解耦状态空间模型ASD-SSM中关键设计选择对模型性能的影响，重点考察尺度约束因子、参数化策略以及尺度数量配置等因素对语义分割性能的作用机制。所有实验均在S3DIS数据集上进行，且均以加载GGAM模块后的模型作为基线，以保证结果具有可比性。

首先分析尺度约束因子对模型性能的影响。尺度约束因子直接控制状态转移矩阵取值范围，从而决定状态在时间维度上的衰减速度。较小的尺度约束因子使状态更快衰减，有利于模型对局部几何变化做出快速响应；较大的尺度约束因子则能够延缓状态衰减，从而增强对长程依赖和全局语义的建模能力。\cref{tab:alpha_ablation}给出了不同尺度约束因子配置下的实验结果。

\begin{table}[htbp!]
	\centering
	\caption{尺度约束因子$\alpha_s$消融实验（3尺度配置）}
	\label{tab:alpha_ablation}
	\begin{tabular}{@{}llll@{}}
		\toprule
		$\alpha_1$ & $\alpha_2$ & $\alpha_3$ & mIoU \\
		\midrule
		0.3 & 0.3 & 0.3 & 73.7 \\
		0.7 & 0.7 & 0.7 & 73.5 \\
		0.5 & 0.5 & 0.5 & 74.0 \\
		0.3 & 0.5 & 0.7 & 74.7 \\
		\textbf{0.3} & \textbf{0.6} & \textbf{0.9} & \textbf{75.2} \\
		0.3 & 0.7 & 0.9 & 75.0 \\
		0.2 & 0.6 & 0.9 & 74.6 \\
		0.4 & 0.6 & 0.9 & 74.7 \\
		0.9 & 0.6 & 0.3 & 67.9 \\
		\bottomrule
	\end{tabular}
\end{table}

实验结果表明，当所有尺度使用相同的尺度约束因子时，模型性能始终低于差异化配置。这说明不同尺度在建模目标上存在本质差异，统一的状态衰减特性难以同时满足局部细节捕捉和全局语义建模的需求。当尺度约束因子从细尺度到粗尺度逐级增大时，模型性能显著提升，其中0.3、0.6、0.9的递增配置取得了最优结果。这一现象验证了细尺度快速衰减、粗尺度缓慢衰减的设计思想，即细尺度更关注局部结构变化，而粗尺度负责维持稳定的长程记忆。相反，当尺度约束因子采用反向配置时，模型性能下降至67.9，说明在细尺度引入过强的长程记忆会干扰局部特征建模，而在粗尺度过快衰减则会破坏全局语义一致性。

随后，我们将ASD-SSM与多种常见参数化策略进行了对比分析，结果如\cref{tab:parameterization_comparison}所示。共享参数方案中三个尺度共享同一组$\overline{A}$矩阵，无法适应不同尺度的差异化建模需求，性能最低。独立参数方案为每个尺度分配专属的$\overline{A}_s$，能够捕获尺度特定的特性，但在每个尺度内所有Patch共享同一组固定参数。参数插值方案通过参数生成器为最细和最粗尺度动态生成$\overline{A}_s$，中间尺度则通过$\overline{A}_1$和$\overline{A}_S$之间的线性插值获得；然而线性插值难以刻画跨尺度间的非线性关系，性能提升有限。相比之下，ASD-SSM优于独立参数方案，表明动态的尺度感知参数生成是比静态逐尺度分配更为有效的参数化策略。

\begin{table}[htbp!]
	\centering
	\caption{不同参数化方法的性能对比}
	\label{tab:parameterization_comparison}
	\begin{tabular}{@{}ll@{}}
		\toprule
		参数化方法 & mIoU \\
		\midrule
		共享参数 & 71.9 \\
		独立参数 & 74.0 \\
		参数插值 & 72.1 \\
		\textbf{ASD-SSM} & \textbf{75.2} \\
		\bottomrule
	\end{tabular}
\end{table}



最后，我们对尺度数量及其对应的窗口扩展配置进行了系统研究，结果如\cref{tab:scale_number}所示。实验结果表明，在单尺度基础上引入第二个尺度即可带来一致的性能提升，其中双尺度最优配置达到73.8\%。进一步引入第三个尺度后性能再次提升，最优三尺度配置达到75.2\%，说明局部、中间和全局三个层次能够提供更充分的结构表示。在窗口扩展策略方面，适度的扩展因子有助于增强上下文建模能力，而过大的窗口反而会稀释局部几何细节。尽管四尺度配置可以带来0.2\%的微弱提升，但训练时间显著增加且收益递减。因此，三尺度配置（扩展因子$(1, 2, 4)$）在性能与效率之间达到了最优平衡。

\begin{table}[htbp!]
	\centering
	\caption{尺度数量与窗口配置消融实验}
	\label{tab:scale_number}
	\begin{tabular}{@{}llll@{}}
		\toprule
		$S$ & $F$ & 训练时间 & mIoU \\
		\midrule
		1 & (1) & 35h & 72.5 \\
		2 & (1,2) & 39h & 73.5 \\
		2 & (1,3) & 40h & 73.8 \\
		2 & (1,4) & 40h & 73.1 \\
		3 & (1,2,2) & 45h & 75.0 \\
		3 & (1,2,3) & 46h & 75.0 \\
		\textbf{3} & \textbf{(1,2,4)} & \textbf{46h} & \textbf{75.2} \\
		3 & (1,3,3) & 45h & 74.8 \\
		4 & (1,2,2,2) & 53h & 75.2 \\
		4 & (1,2,2,3) & 54h & 75.4 \\
		\bottomrule
	\end{tabular}
\end{table}

\section{基于图多项式展开的多尺度建模：Chebyshev多尺度SSM}

前述ASD-SSM在序列域通过不同窗口大小实现多尺度特征学习，其感受野相对固定，难以在同一模块中同时建模不同空间尺度的几何信息。基于图拉普拉斯的多项式建模为解决上述问题提供了新的视角。通过Chebyshev多项式递推，可以将点云特征分解为多个空间尺度的分量：低阶分量感受野小，捕获点自身及直接邻域的局部几何差异；高阶分量感受野更大，聚合更广范围的上下文结构信息。这种多尺度分解具有三方面优势：首先，能够针对不同空间尺度采用差异化建模策略，避免统一处理导致的建模冲突；其次，各尺度分量的语义纯净度更高，避免了局部细节与全局结构混合建模导致的相互干扰；第三，与序列域方法具有天然互补性，两者结合能够更全面地刻画点云的多尺度特征。

为此，本节提出基于Chebyshev多项式近似的多尺度图状态空间模型（Chebyshev多尺度SSM），与前述ASD-SSM形成互补。Chebyshev多尺度SSM的核心思想是：利用Chebyshev多项式对图拉普拉斯算子进行递推展开，将点云特征分解为多个空间尺度的分量，然后使用独立的Mamba模块分别对各分量进行序列建模，最后通过可学习的加权融合得到增强后的特征表示。这种"多尺度分解 + 分量独立建模"的设计，使模型能够针对不同空间尺度采用差异化的建模策略，从而更有效地捕获点云的多尺度几何信息。Chebyshev多尺度SSM通过Chebyshev多项式递推计算，避免了计算密集的特征分解，将复杂度从$O(N^3)$降低至$O(KN)$（其中$K$为多项式阶数，通常为2-5）；通过为每个空间尺度分量配备独立的Mamba模块，使模型能够学习适应不同尺度特性的建模策略；通过可学习的融合权重，模型能够自适应地平衡不同尺度分量的贡献。整体流程包括四个核心模块：序列化窗口图构建、Chebyshev多项式多尺度展开、分量独立建模和多尺度融合，整体架构如\cref{fig:chebyshev_arch}所示。

\begin{figure}[htbp]
	\centering
	\includegraphics[width=\linewidth]{pinyu.jpg}
	\caption{Chebyshev多尺度SSM架构图}
	\label{fig:chebyshev_arch}
\end{figure}

\subsection{序列化窗口图构建}

多尺度分解需要首先构建点云的邻域关系。基于序列化的邻域构建方法可以将复杂度从$O(N\log N)$降至$O(N)$。与GGAM采用的窗口划分策略不同，本节采用序列邻居策略，更适合图多项式递推所需的稀疏连接。

具体地，我们利用Point Transformer V3已有的空间序列化顺序，直接使用序列邻域构建图：
\begin{equation}
\mathcal{E} = \{(i, j) \mid q_j \in \mathcal{N}(q_i)\}
\end{equation}
其中$\mathcal{N}(q_i) = \{q_{i-P/2}, \ldots, q_{i-1}, q_{i+1}, \ldots, q_{i+P/2}\}$为点$q_i$在序列上前后各$P/2$个点，$P$表示局部邻域大小。对于序列边界点，采用截断策略处理邻域不足的情况。边权重采用高斯核：
\begin{equation}
w_{ij} = \exp\left(-\frac{\|\mathbf{q}_i - \mathbf{q}_j\|^2}{2\sigma^2}\right)
\end{equation}
其中$\sigma$为尺度参数，设为当前批次所有边长度的均值。该权重设计使得空间距离近的点具有更强的连接权重，距离远的点影响快速衰减，为后续多尺度分解提供合理的邻域拓扑结构。对于包含多个场景的批次数据，我们为每个场景独立构建图，避免跨场景的信息泄漏。由于序列化顺序已预先计算，无需额外排序开销。

\subsection{Chebyshev多项式多尺度分解}

获得邻域边集$\mathcal{E}$后，对点云特征进行多尺度分解。设输入特征为$X \in \mathbb{R}^{N \times C}$，我们使用Chebyshev多项式计算$K$个尺度分量的特征表示。

基于边权重$w_{ij}$，构造邻接权重矩阵$\mathbf{W} \in \mathbb{R}^{N \times N}$，其中$\mathbf{W}_{ij} = w_{ij}$若$(i,j) \in \mathcal{E}$，否则为0。度矩阵$D \in \mathbb{R}^{N \times N}$为对角矩阵，对角元素$D_{ii} = \sum_{q_j \in \mathcal{N}(q_i)} w_{ij}$表示点$q_i$的加权度。为适配Chebyshev多项式对输入特征值范围$[-1, 1]$的要求，需要对归一化拉普拉斯$L_{norm} = I - D^{-1/2}\mathbf{W}D^{-1/2}$进行线性变换。由于$L_{norm}$的特征值范围为$[0, 2]$，构造$\tilde{L} = L_{norm} - I$可将特征值映射到$[-1, 1]$，代入$L_{norm}$的定义得到：
\begin{equation}
\tilde{L} = L_{norm} - I = -D^{-1/2}\mathbf{W}D^{-1/2}
\end{equation}

对于输入特征矩阵$X \in \mathbb{R}^{N \times C}$，我们递推计算$K$个Chebyshev基作用的结果：
\begin{equation}
\begin{aligned}
T_0(\tilde{L})X &= X \\
T_1(\tilde{L})X &= \tilde{L}X \\
T_\ell(\tilde{L})X &= 2\tilde{L}T_{\ell-1}(\tilde{L})X - T_{\ell-2}(\tilde{L})X, \quad \ell = 2, \ldots, K-1
\end{aligned}
\end{equation}
关键计算是$\tilde{L}X$，利用$\tilde{L} = -D^{-1/2}\mathbf{W}D^{-1/2}$的稀疏性，该操作可高效实现：
\begin{equation}
(\tilde{L}X)_i = -\sum_{q_j \in \mathcal{N}(q_i)} \frac{w_{ij}}{\sqrt{D_{ii}D_{jj}}} X_j
\end{equation}
由于使用scatter操作向量化计算，其复杂度为线性。不同阶的Chebyshev多项式对应不同的空间感受野：$T_0(\tilde{L})X = X$保留点的原始特征（0跳）；$T_1(\tilde{L})X = \tilde{L}X$聚合1跳邻域内的加权差异信息；$T_\ell(\tilde{L})X$的感受野随阶数$\ell$增大而扩展，高阶分量能够捕获更大范围的结构信息。这种多尺度分解为后续的差异化建模奠定了基础。

\subsection{多尺度分量独立建模}

获得$K$个空间尺度的分量$\{T_\ell(\tilde{L})X\}_{\ell=0}^{K-1}$后，本文为每个分量配备独立的Mamba模块，实现尺度特定的序列建模。这一设计源于不同空间尺度分量的固有特性差异：对于低阶分量（$\ell=0,1$），感受野小，捕获点自身及直接邻域的局部几何信息，局部差异和边界细节更为关键；而对于高阶分量（$\ell \geq 2$），感受野大，聚合更广范围的上下文信息，长程依赖关系更为重要。传统方法使用单一Mamba处理所有尺度的混合特征，难以同时适应小感受野分量的局部特性和大感受野分量的全局特性，而独立Mamba使模型能够学习特定尺度的最优建模策略。

对于第$\ell$个分量的特征$X^{(\ell)} = T_\ell(\tilde{L})X \in \mathbb{R}^{N \times C}$，沿用Point Transformer V3已有的序列化顺序$Q = (q_1, q_2, \ldots, q_N)$，将其按该顺序排列为序列后送入Mamba进行序列建模：
\begin{equation}
X^{(\ell)}_{seq} = \left(X^{(\ell)}_{q_1},\, X^{(\ell)}_{q_2},\, \ldots,\, X^{(\ell)}_{q_N}\right)
\end{equation}
第$\ell$个分量的Mamba模块定义为：
\begin{equation}
Y^{(\ell)} = \text{Mamba}_\ell(X^{(\ell)}_{seq})
\end{equation}
其中$\text{Mamba}_\ell$为第$\ell$个分量专属的Mamba块，包含双向SSM（前向+后向）和残差连接。各分量Mamba的参数完全独立，$\text{Mamba}_\ell$与$\text{Mamba}_j$（$\ell \neq j$）不共享权重。每个分量输出后应用LayerNorm以稳定训练，Mamba输出按序列化的逆序恢复为原始点云顺序，得到$Y^{(\ell)} \in \mathbb{R}^{N \times C}$。

\subsection{多尺度特征融合}

获得$K$个尺度分量的Mamba输出$\{Y^{(\ell)}\}_{\ell=0}^{K-1}$后，采用可学习加权融合策略将其融合为统一的特征表示。定义可学习的融合权重$\boldsymbol{\beta} = [\beta_0, \beta_1, \ldots, \beta_{K-1}] \in \mathbb{R}^K$，初始化为均匀权重$\beta_\ell = 1/K$，融合特征计算为：
\begin{equation}
Y_{multi} = \sum_{\ell=0}^{K-1} \alpha_\ell Y^{(\ell)}, \quad \alpha_\ell = \frac{\exp(\beta_\ell)}{\sum_{r=0}^{K-1}\exp(\beta_r)}
\end{equation}
其中$\alpha_\ell$为归一化后的权重（通过softmax保证$\sum_\ell \alpha_\ell = 1$）。最终输出通过残差连接和MLP层得到：
\begin{equation}
\begin{aligned}
Z &= \text{Concat}(X, Y_{multi}) \\
\hat{X} &= \text{MLP}(Z) \\
X' &= X + \hat{X}
\end{aligned}
\end{equation}
其中$\text{Concat}(X, Y_{multi}) \in \mathbb{R}^{N \times 2C}$将原始特征和多尺度特征拼接，$\text{MLP}$为2层感知机（隐藏层维度为$C$，使用GELU激活），输出维度为$C$。拼接操作保留了原始特征信息，避免多项式变换导致的信息损失；残差连接使模型能够自适应地选择多尺度增强的强度，在不同编码层级和数据集上具有更好的适应性。

\subsection{消融实验}

通过系统消融实验，分析Chebyshev多尺度SSM中关键设计选择对模型性能的影响，重点考察多项式阶数和Mamba独立性等因素。所有实验均在S3DIS数据集上进行，且均在仅启用Chebyshev多尺度SSM、不加载GGAM和ASD-SSM的条件下进行，以隔离各设计选择的独立影响。

\textbf{Chebyshev多项式阶数。}Chebyshev多项式阶数$K$决定了多尺度分解的精细度。表\ref{tab:cheb_order}展示了不同$K$值的性能对比。

\begin{table}[htbp!]
\centering
\caption{Chebyshev多项式阶数消融实验}
\label{tab:cheb_order}
\begin{tabular}{@{}cccc@{}}
\toprule
阶数$K$ & 参数量 & 推理时间 (ms) & mIoU (\%) \\
\midrule
1 & 1.00$\times$ & 103.9 & 70.5 \\
2 & 1.76$\times$ & 140.9 & 70.7 \\
\textbf{3} & \textbf{2.52$\times$} & \textbf{201.8} & \textbf{70.8} \\
4 & 3.28$\times$ & 250.8 & 70.8 \\
5 & 4.04$\times$ & 306.9 & 71.0 \\
\bottomrule
\end{tabular}
\end{table}

实验结果表明，$K=1$时仅利用原始特征与一跳邻域信息，性能最低，说明单一空间尺度不足以捕获复杂的几何结构。随着阶数从$K=1$增加到$K=3$，性能稳步提升，验证了多尺度谱分解的有效性。当$K$继续增大至4和5时，性能趋于饱和，而参数量与推理时间显著增长，边际收益远不足以抵消额外的计算开销。综合性能与效率，本文选取$K=3$作为最佳平衡点。

\textbf{Mamba独立性。}本节的核心创新是为每个尺度分量配备独立的Mamba。表\ref{tab:mamba_sharing}对比了不同参数共享策略的效果。

\begin{table}[htbp!]
\centering
\caption{Mamba参数共享策略对比}
\label{tab:mamba_sharing}
\begin{tabular}{@{}lcc@{}}
\toprule
参数共享策略 & 参数量 & mIoU (\%) \\
\midrule
所有分量共享Mamba & 1.00$\times$ & 70.4 \\
低阶/高阶各自共享（2组） & 1.59$\times$ & 70.7 \\
\textbf{每个分量独立Mamba} & \textbf{2.18$\times$} & \textbf{70.8} \\
\bottomrule
\end{tabular}
\end{table}

实验显示，三种策略的推理时间几乎相同（约200 ms），说明计算瓶颈在于图构建和Chebyshev基计算，Mamba本身并非性能瓶颈。因此，独立Mamba仅增加参数量而不增加推理开销。所有分量共享单一Mamba时，参数量最少但难以同时适应不同空间尺度分量的特性差异。将低阶分量（$\ell=0,1$）和高阶分量（$\ell=2$）分为两组各自共享Mamba，性能有所提升，验证了分组建模的价值。每个分量使用独立Mamba时达到最优性能，虽然参数量增至2.18倍，但由于不影响推理速度，这一开销是可接受的。

\section{本章小结}

本章从序列域和谱域两个互补视角构建了完整的多尺度点云状态空间建模框架。在序列域，ASD-SSM通过动态参数生成机制解决了状态转移参数输入无关的局限性，使模型能够根据每个Patch的局部几何内容自适应调整$\overline{A}$。通过尺度约束因子显式控制不同尺度的状态衰减特性，细尺度快速响应局部几何细节，粗尺度维持长程记忆建模全局语义，在S3DIS数据集上达到75.2\% mIoU。在谱域，Chebyshev多尺度SSM利用Chebyshev多项式递推计算实现$O(KN)$复杂度的多尺度分解，将点云特征展开为$K$个不同感受野的分量，通过为每个分量配备独立Mamba模块实现尺度特定的序列建模，并通过可学习权重自适应融合多尺度特征，在S3DIS数据集上达到70.8\% mIoU。两种方法在多尺度建模上形成正交互补：ASD-SSM在序列化空间中以不同窗口大小构建层次化依赖，Chebyshev多尺度SSM通过图谱分解在不同空间感受野上提取几何信息。两者叠加使用具有协同效应，整体实验结果将在第六章进行汇总分析。

\chapter{实验}
\textbf{Implementation Details:}
我们在室内数据集S3DIS和分类数据集ModelNet40上进行了点云分析的基准测试，所有实验都在Nvidia A6000 GPU上进行。网络训练中，我们使用 AdamW 优化器及 OneCycleLR 学习率调度策略，初始最大学习率设为 6e-3，并采用前期预热及后期余弦退火循环的调度机制。整个训练过程持续 3000 轮（epoch），批量大小为 6，并启用混合精度（AMP）加速以提高训练效率。在编码器中，channel大小设置为(32, 64, 128, 256, 512)。解码器中，channel大小设置为(64, 64, 128, 256)。Patch Size统一设置为128,与现有工作保持一致以利于
公平对比。
我们选择交叉熵损失函数作为损失判断依据。交叉熵损失是最常用的分类以及分割任务损失函数，广泛应用于语义分析问题。

PointSS采用从头训练的方式，不使用任何预训练策略，使得性能提升可归因于架构设计本身，而非预训练策略的差异。


\section{基准数据集与评估指标}
\textbf{S3DIS:}
S3DIS 数据集是由斯坦福大学发布的，用于室内场景理解和点云处理的标准数据集之一。它包含了多个室内建筑的 3D 点云数据，主要用于点云语义分割任务。其包含6种大型建筑，涵盖了办公楼、休息区、会议室、大厅等多种室内环境。S3DIS包含 272 个房间，每个房间的 3D 点云数据均经过精细标注。数据集中每个点包含三维坐标和RGB颜色信息，同时包含每个点的语义类别。数据集共标注了13 个类别：如墙（wall）、地板（floor）、桌子（table）、椅子（chair）、窗户（window）、书架（bookshelf） 等。室内的各种环境被平均划分到了六个区域之中，我们选择区域五进行结果验证。


\textbf{ModelNet40:}
ModelNet40 是由 Princeton University 提供，旨在为 3D 对象分类和识别任务提供标准化的基准。其包括40 个物体类别，12,311 个 3D 对象样本。训练集包含8,156 个样本，测试集包含4,155 个样本。数据集中的每个对象都经过详细的分类，涵盖各种日常物品，如椅子、桌子、瓶子、飞机、汽车等。


\textbf{mIoU:}mean Intersection over Union是语义分割等任务中常用的评估指标，用于衡量预测结果与真实标签之间的重叠程度。对每一个类别$c$，IoU（交并比）定义为：\begin{equation}\mathrm{IoU}_c=\frac{\mathrm{TP}_c}{\mathrm{TP}_c+\mathrm{FP}_c+\mathrm{FN}_c}\end{equation}
$\mathrm{TP}_{c}$为真阳性，$\mathrm{FP}_{c}$为假阳性，$\mathrm{FN}_{c}$为假阴性。取所有类别的平均值，得到mIoU。其对类不均衡的数据更鲁棒，比Accuracy更常用于3D语义分割任务。因此我们选择mIoU为指标来评价在S3DIS上的点云语义分割任务。

\textbf{OA:}Overall Accuracy，即模型在所有点云样本上预测正确的数量除以所有样本的总数。其简单直观，常用于分类任务的总体指标，因此我们选择OA为指标来评价在ModelNet40上的点云分类任务。

\section{定量实验}

\subsubsection{Experiment results on the S3DIS dataset}
\cref{tab:s3dis}汇总了S3DIS数据集上的语义分割性能对比。在不使用预训练的条件下，PointSS在Area 5上达到75.9\%的mIoU，优于所有对比方法。具体而言，相比当前最优的Point Transformer V3提升了2.7个百分点，相比基于SSM的PCM提升了6.1个百分点。

\begin{table}[htbp!]
	\caption{在S3DIS上的分割性能对比}
	\label{tab:s3dis}
	\centering
	\begin{tabular}{@{}llll@{}}
		\toprule
		\textbf{模型}           & \textbf{出处} & \textbf{架构} & \textbf{mIoU} \\ \midrule
		Point Transformer V2\cite{ptv2}  & NeurIPS'22 & Transformer & 72.6          \\
		PointMetaBase\cite{pmb}         & CVPR'23 & Transformer & 72.0          \\
		SuperpointTransformer\cite{spt} & ICCV'23 & Transformer & 68.1          \\
		Swin3D\cite{Swin3D}                & CVMJ'23 & Transformer & 72.5          \\
		PPT+SpareUNET\cite{ppt}         & CVPR'24 & Transformer & 72.7          \\
		OA-CNNs\cite{oacnn}               & CVPR'24 & CNN         & 71.1          \\
		OneFormer3D\cite{OneFormer3D}           & CVPR'24 & Transformer & 72.4          \\
		Point Transformer V3\cite{ptv3}           & CVPR'24 & Transformer & 73.2          \\
		PCM\cite{pcm}                   & AAAI'25 & SSM         & 69.8         \\
		PointMSGT\cite{pointmsgt}      & Mathematics'24 & Transformer & 68.6         \\
		PointGA\cite{pointga}           & Sci. Rep.'25 & Transformer & 66.2         \\
		\textbf{PointSS}               & & \textbf{SSM}           & \textbf{75.9}          \\ \bottomrule
	\end{tabular}
	
	
\end{table}


\subsubsection{Experiment results on the ModelNet40 dataset}

\cref{tab:ModelNet}汇总了ModelNet40数据集上的分类性能对比。PointSS达到94.6\%的整体准确率（OA），超越了所有基于SSM的对比方法（包括Mamba3D和PointMamba），同时与基于Transformer的方法保持竞争力。值得注意的是，在该接近饱和的基准上，大多数方法的准确率集中在较窄的范围内，此量级的边际增益不应被过度解读。

\begin{table}[htbp!]
	\centering
	\caption{在ModelNet40上的分类性能对比}
	\label{tab:ModelNet}
	\begin{tabular}{@{}llll@{}}
		\toprule
		\textbf{模型}           & \textbf{出处} & \textbf{架构} & \textbf{OA} \\ \midrule
		Point Transformer V2\cite{ptv2}  & NeurIPS'22 & Transformer & 94.2          \\
		Point-FEMAE\cite{FEMAE}        & AAAI'24 & Transformer & 94.5             \\
		Point2Vec\cite{Point2Vec}      & LNCS'24 & Transformer & 94.8          \\
		Mamba3D\cite{Mamba3D}              & ACM MM'24 & SSM & 94.1          \\
		PointMamba\cite{PointMamba}   & NeurIPS'24 & SSM         & 93.6         \\
		OTMae3D\cite{OTMae3D}    & TIP'25 & Transformer & 94.5          \\
		PCM\cite{pcm}     & AAAI'25 & SSM         & 93.4         \\
		PointGA\cite{pointga}           & Sci. Rep.'25 & Transformer & 93.8         \\
		\textbf{PointSS}               & & \textbf{SSM}             & \textbf{94.6}          \\ \bottomrule
	\end{tabular}
\end{table}


\subsubsection{Ablation study}
如\cref{tab:progressive_improvement}所示，仅在基线模型中引入GGAM即可带来2.2\%的mIoU提升。这一结果表明，通过窗口化邻域提取几何先验并结合双序列化特征融合，GGAM能够有效缓解点云序列化过程中空间邻近点被割裂所导致的几何信息丢失问题，使模型在特征更新阶段能够整合更加完整的局部几何结构信息，从而显著提升语义分割性能。

在此基础上，进一步引入ASD-SSM后，模型性能从72.5\%提升至75.2\%，带来2.7\%的mIoU增益。该结果表明，通过尺度感知参数生成器为不同尺度动态生成定制化的状态空间参数，ASD-SSM能够使细尺度聚焦于局部几何细节，粗尺度建模长程语义上下文，从而有效提升模型的多尺度表达能力。

综合来看，GGAM与ASD-SSM在功能上具有明显的互补性：GGAM主要解决序列化引入的几何信息缺失问题，为模型提供更加可靠的局部几何感知基础；ASD-SSM则通过参数解耦的多尺度建模机制，实现了粗尺度长程语义建模与细尺度局部细节响应之间的有效平衡。两者协同作用，使得PointSS在保持计算效率的同时，实现了显著的整体性能提升。

在空域方法（GGAM + ASD-SSM）的基础上，进一步引入Chebyshev多尺度SSM后，模型性能从75.2\%提升至75.9\%（+0.7\%）。Chebyshev多尺度SSM通过图拉普拉斯递推聚合不同跳数的邻域信息，提供了空域方法未覆盖的拓扑结构特征，与空域方法从不同维度刻画点云的多尺度几何特征，最终使完整系统总提升幅度达到5.6\% mIoU。

\begin{table}[htbp!]
	\centering
	\caption{模块渐进式性能提升路径（S3DIS数据集）}
	\label{tab:progressive_improvement}
	\begin{tabular}{@{}lll@{}}
		\toprule
		\textbf{模型配置} & \textbf{mIoU (\%)} & \textbf{提升} \\
		\midrule
		纯Mamba基线 & 70.3 & - \\
		+ GGAM & 72.5 & +2.2 \\
		+ ASD-SSM & 75.2 & +2.7 \\
		\textbf{+ Chebyshev多尺度SSM} & \textbf{75.9} & \textbf{+0.7} \\
		\midrule
		\multicolumn{2}{l}{\textbf{总提升}} & \textbf{+5.6} \\
		\bottomrule
	\end{tabular}
\end{table}



\section{计算效率分析}

为验证状态空间模型线性复杂度的计算优势，我们在相同条件下（FP32精度、batch size为1、NVIDIA A6000 GPU）比较了PointSS与Point Transformer V3的推理时间，结果如\cref{fig:efficiency}所示。

\begin{figure}[htbp!]
	\centering
	\begin{tikzpicture}
	\begin{axis}[
		width=0.85\linewidth,
		height=6cm,
		xlabel={Patch Size},
		ylabel={推理时间 (ms)},
		xtick={128,256,512,1024},
		xticklabels={128,256,512,1024},
		xmin=64, xmax=1150,
		ymin=0.6, ymax=1.55,
		legend pos=south east,
		legend style={font=\small},
		grid=major,
		grid style={dashed,gray!30},
	]

	\addplot[color=blue, mark=square*, thick, solid] coordinates {
		(128, 1.089)
		(256, 1.147)
		(512, 1.274)
	};
	\addlegendentry{PTv3}

	\addplot[color=orange, mark=*, thick, solid] coordinates {
		(128,  0.950)
		(256,  1.038)
		(512,  1.032)
		(1024, 1.308)
	};
	\addlegendentry{PointSS (Ours)}

	\draw[color=blue, dashed, thick] (axis cs:512,1.274) -- (axis cs:1024,1.274);
	\node[color=blue, font=\small\bfseries, anchor=south] at (axis cs:1024,1.274) {OOM};

	\end{axis}
	\end{tikzpicture}
	\caption{PTv3与PointSS在不同Patch Size下的推理时间对比（FP32精度、batch size为1、NVIDIA A6000 GPU）}
	\label{fig:efficiency}
\end{figure}

在Patch Size从128到512的范围内，PointSS的推理时间相比PTv3降低约9\%--19\%。当Patch Size增至1024时，PTv3因自注意力机制的二次方内存开销导致显存溢出（OOM），而PointSS仍可正常运行，体现了线性复杂度在大规模点云处理中的可扩展性优势。


\section{定性实验}
\cref{fig:vis}展示了PointSS、Point Transformer和Point Cloud Mamba在S3DIS数据集上四个场景中的分割结果对比。如图中圆圈标注所示，PointSS在以下三类区域中表现出优势。

对于窗户等具有较大空间跨度的结构（office\_10上方圆圈），基线方法产生了明显的碎片化或不完整的预测结果；而PointSS生成了更完整的分割区域，与真实标签具有更高的一致性。这一改进归因于GGAM对窗口化邻域内几何关系的显式建模，补偿了序列化过程中空间邻近性丢失所导致的上下文不连续问题。

在包含杂物等小型不规则物体的区域（office\_12和office\_14下方圆圈），基线方法容易将这些区域与周围的桌椅混淆，导致边界预测不稳定。PointSS利用多尺度架构中细尺度的快速状态衰减，能够更有效地响应局部几何变化，在此类区域获得更准确的类别预测。

在书架与墙面之间的类别边界处（office\_13），基线方法出现了模糊或错误分类的预测，而PointSS产生了更清晰、更准确的边界。通过ASD-SSM粗尺度中较慢的状态衰减，模型保留了更广泛的上下文信息，从而提高了边界预测的一致性。

上述定性结果与定量实验中观察到的趋势一致，进一步验证了PointSS在复杂室内场景中的有效性。
\begin{figure}[htbp]
	\centering
	\includegraphics[width=\linewidth]{visualize.JPG}
	\caption{S3DIS Area 5上四个代表性室内场景的语义分割结果定性对比}
	\label{fig:vis}
\end{figure}



\chapter{结论}
本文针对现有基于Mamba的点云分析方法存在的序列化几何信息丢失与多尺度建模能力不足两方面问题，提出了PointSS方法。PointSS通过三个核心模块——全局几何感知机制（GGAM）、自适应尺度解耦状态空间模型（ASD-SSM）和基于Chebyshev多项式近似的多尺度图状态空间模型（Chebyshev多尺度SSM），分别从空域几何补偿、多尺度参数解耦和图多项式拓扑建模三个维度提升了点云分析性能。

在几何感知方面，GGAM通过窗口化邻域提取曲率和法向量等几何先验，结合双序列化融合与跨序列注意力机制，为Mamba显式注入几何信息，有效补偿了序列化过程中空间邻近性丢失所导致的几何学习困难。消融实验表明，GGAM对序列化割裂严重的区域和几何复杂的类别带来了显著的性能提升。

在多尺度特征学习方面，ASD-SSM通过尺度感知参数生成器为不同尺度动态生成定制化的状态转移参数，使细尺度快速响应局部几何细节变化，粗尺度保持长程记忆以捕获全局语义上下文，实现了层次化的多尺度特征学习。实验表明，相比共享参数或完全独立参数的方案，ASD-SSM在保持参数高效的前提下取得了更优的性能。

在图多项式建模方面，Chebyshev多尺度SSM利用Chebyshev多项式对图拉普拉斯算子的递推展开，以线性复杂度将点云特征分解为多个不同感受野的分量，并为每个分量配备独立的Mamba模块进行差异化序列建模。该方法从图多项式的视角补偿了空域方法在拓扑信息建模上的不足，与GGAM和ASD-SSM形成互补。实验表明，Chebyshev多尺度SSM与空域方法结合后产生了协同增益，验证了两类方法的互补性。

在S3DIS和ModelNet40标准基准上的实验表明，PointSS在不依赖预训练的条件下，在语义分割任务上超越了现有的SSM方法和Transformer方法，在分类任务上超越所有SSM类对比方法并与Transformer方法保持竞争力。计算效率分析表明，PointSS在保持线性时间复杂度的同时具有推理速度优势，并在大Patch Size场景下展现出更好的可扩展性。定性实验进一步验证了PointSS在大空间跨度结构、细粒度物体和类别边界等场景中的有效性。

尽管取得了上述成果，PointSS仍存在一些局限性。尺度约束因子$\alpha_s$作为固定超参数需要手动设定，在推广到新数据集或不同点云分布时可能需要针对性调优。此外，GGAM中的双序列化设计相较于单流方法引入了额外的计算开销。目前PointSS仅在室内场景分割和合成物体分类基准上进行了验证，其在室外LiDAR数据集上的有效性有待进一步检验。

未来工作将探索自适应学习尺度约束因子以减少人工调参的需求，同时将PointSS扩展到大规模室外数据集，并研究更高效的几何特征融合策略，以应对更复杂的实际应用场景。
%% The Appendices part is started with the command \appendix;
%% appendix sections are then done as normal sections

%% bibitems, please use
%%
%%  \bibliographystyle{elsarticle-num} 
%%  \bibliography{<your bibdatabase>}

%% else use the following coding to input the bibitems directly in the
%% TeX file.

%% Refer following link for more details about bibliography and citations.
%% https://en.wikibooks.org/wiki/LaTeX/Bibliography_Management

% \begin{thebibliography}{00}

% %% For numbered reference style
% %% \bibitem{label}
% %% Text of bibliographic item

% \bibitem{lamport94}
%   Leslie Lamport,
%   \textit{\LaTeX: a document preparation system},
%   Addison Wesley, Massachusetts,
%   2nd edition,
%   1994.

% \end{thebibliography}
\bibliographystyle{unsrt}
\bibliography{bishe_ref}
\end{document}

\endinput
%%
%% End of file `elsarticle-template-num.tex'.
